\appendixitem{FORESTS AND TREES}

Let \(\Gamma = (\mathrm{A},\mathrm{S})\) be a graph. A circuit of \(\Gamma\) is a sequence

\[
(x _ {1}, \ldots , x _ {n})
\]

of distinct vertices of \(\Gamma\) such that \(n \geqslant 3\) , \(x_{i}\) is joined to \(x_{i+1}\) for \(1 \leqslant i < n\) and \(x_{n}\) is joined to \(x_{1}\) . \(\Gamma\) is called a forest if there is no circuit in \(\Gamma\) ; every subgraph of \(\Gamma\) is then also a forest. A connected forest is called a tree; the connected components of a forest are thus trees.

PROPOSITION 2. Let \(\Gamma = (\mathrm{A},\mathrm{S})\) be a forest having only a finite number of vertices.

\begin{enumerate}
\def\labelenumi{(\roman{enumi})}
\item
  If \(\Gamma\) has at least one vertex, it has a terminal vertex.
\item
  If \(\Gamma\) has at least two vertices, there is a partition \((\mathbf{S}', \mathbf{S}'')\) of its set of vertices into two non-empty subsets such that two distinct vertices that both belong to \(\mathbf{S}'\) or both belong to \(\mathbf{S}''\) are never joined.
\end{enumerate}

Suppose that \(\Gamma\) has at least one vertex and let \((x_{0},\ldots,x_{n})\) be an injective path of maximal length in \(\Gamma\) . The vertex \(x_{0}\) cannot be joined to a vertex y distinct from \(x_{0},x_{1},\ldots,x_{n}\) , since otherwise there would exist an injective path in \(\Gamma\) of length \(n+1\) , namely \((y,x_{0},\ldots,x_{n})\) . The vertex \(x_{0}\) is not joined to any vertex \(x_{i}\) with \(2\leqslant i\leqslant n\) , since otherwise \((x_{0},x_{1},\ldots,x_{i})\) would be a circuit in the forest \(\Gamma\) . Thus, \(x_{0}\) is terminal.

We shall prove (ii) by induction of the number \(m\) of vertices of \(\Gamma\) , the case \(m = 2\) being trivial. Suppose then that \(m \geqslant 3\) and that assertion (ii) is proved for graphs with \(m - 1\) vertices. Let \(a\) be a terminal vertex of \(\Gamma\) (cf.~(i)). We apply the induction hypothesis to the full subgraph of \(\Gamma\) whose vertices are the vertices \(x \neq a\) of \(\Gamma\) . Thus, there exist two non-empty disjoint subsets \(S_1'\) and \(S_1''\) of \(S\) with \(S_1' \cup S_1'' = S - \{a\}\) , and such that two distinct vertices in \(S_1'\) (resp. \(S_1''\) ) are never joined. Since \(a\) is joined to at most one vertex of \(\Gamma\) , we can suppose for example that it is not joined to any vertex in \(S_1'\) . The partition \((S_1', S_1'' \cup \{a\})\) then has the required property. Q.E.D.

For any integer \(n \geqslant 1\) , denote by \(A_{n}\) the graph whose vertices are the integers \(1, 2, \ldots, n\) and whose edges are the pairs \(\{i, j\}\) with \(i - j = \pm 1\) :

1

2

3

n---1

n

○

○

○

\ldots\ldots{}

○

A graph \(\Gamma\) is said to be a chain of length \(m \geqslant 0\) if it is isomorphic to \(A_{m+1}\) . This is equivalent to the existence in \(\Gamma\) of an injective path \((x_{0}, \ldots, x_{m})\) containing all the vertices, the vertices \(x_{i}\) and \(x_{j}\) not being joined if \(|i-j| > 1\) .

PROPOSITION 3. A graph is a chain if and only if its number of vertices is finite and non-zero and it is a tree with no ramification point.

Suppose that the graph \(\Gamma\) is a chain \((x_0, \ldots, x_m)\) with the properties listed before the statement of Prop. 3. It is clear that any vertex of \(\Gamma\) is joined to at most two vertices. If \(i < j\) the path \((x_i, \ldots, x_j)\) extracted from the path \((x_0, \ldots, x_m)\) joins \(x_i\) to \(x_j\) ; thus, \(\Gamma\) is connected. Finally, let \((x_{p_1}, \ldots, x_{p_n})\) be a circuit in \(\Gamma\) , and let \(p_k\) be the smallest of the distinct integers \(p_1, \ldots, p_n\) . There exist distinct indices \(i\) and \(j\) such that \(x_{p_k}\) is joined to \(x_{p_i}\) and \(x_{p_j}\) : this follows from the definition of a circuit. Since \(p_k < p_i\) and \(p_k < p_j\) , we necessarily have \(p_i = p_j = p_k + 1\) , which is a contradiction since \(p_1, \ldots, p_n\) are distinct. Thus, there is no circuit in \(\Gamma\) .

Conversely, let \(\Gamma\) be a tree with no ramification point and with a finite non-zero number of vertices, and let \((x_0, \ldots, x_m)\) be an injective path of maximal length in \(\Gamma\) . Denote by T the set of vertices other than \(x_0, \ldots, x_m\) . A vertex \(b \in T\) cannot be joined to any vertex \(x_i\) , for we would have either

\begin{enumerate}
\def\labelenumi{\alph{enumi})}
\item
  \(i = 0\) , but then \((b, x_0, \ldots, x_m)\) would be an injective path of length \(m + 1\) in \(\Gamma\) , or
\item
  \(i = m\) , but then \((x_0, \ldots, x_m, b)\) would be an injective path of length \(m + 1\) in \(\Gamma\) , or
\item
  \(0 < i < m\) , but then \(x_{i}\) would be joined to three distinct vertices \(x_{i-1}\) , \(x_{i+1}\) and \(b\) .
\end{enumerate}

Since \(\Gamma\) is connected, \(T\) is empty by Cor. 1 of Prop. 1. Moreover, if there were \(i,j\) with \(j - i > 1\) and \(x_{i}, x_{j}\) joined, there would be a circuit \((x_{i}, x_{i+1}, \ldots, x_{j})\) in \(\Gamma\) . Thus, \(\Gamma\) is a chain. Q.E.D.

\exercisesection{Exercises for § 1}

\begin{enumerate}
\def\labelenumi{\arabic{enumi})}
\tightlist
\item
  \begin{enumerate}
  \def\labelenumii{\alph{enumii})}
  \tightlist
  \item
    Let (W, S) be a Coxeter system and let \(s_{1}, \ldots, s_{r}\) be elements of S. Put \(w = s_{1} \ldots s_{r}\) . Show that if \(l_{\mathbb{S}}(w) < r\) , there exist two integers p and q with \(1 \leqslant p < q \leqslant r\) such that \(w = s_{1} \ldots s_{p-1}s_{p+1} \ldots s_{q-1}s_{q+1} \ldots s_{r}\) . Show that there is a strictly increasing sequence of integers \(j(1), \ldots, j(k)\) between 1 and r such that \((s_{j(1)}, \ldots, s_{j(k)})\) is a reduced decomposition of w.
  \end{enumerate}
\end{enumerate}

\begin{enumerate}
\def\labelenumi{\alph{enumi})}
\setcounter{enumi}{1}
\tightlist
\item
  Let \((\mathbf{W},\mathbf{S})\) be a Coxeter system and \(X,Y,Z\) three subsets of \(S\) . Show that
\end{enumerate}

\[
W _ {X} \cap (W _ {Y}. W _ {Z}) = (W _ {X} \cap W _ {Y}). (W _ {X} \cap W _ {Z})
\]

(show that every element \(w \in W_{Y}.W_{Z}\) has a reduced decomposition

\[
(s _ {1}, \ldots , s _ {h}, t _ {1}, \ldots , t _ {k})
\]

such that \(s_i \in Y\) and \(t_j \in Z\) and use Cor. 1 of Prop. 7 of no. 8).

Show that

\[
W _ {X}. \left(W _ {Y} \cap W _ {Z}\right) = \left(W _ {X}. W _ {Y}\right) \cap \left(W _ {X}. W _ {Z}\right).
\]

\begin{enumerate}
\def\labelenumi{\arabic{enumi})}
\setcounter{enumi}{1}
\item
  Let \((W,S)\) be a Coxeter system and X a subset of S. Show that \(W_{X}\) is normal in W if and only if every element of X commutes with every element of S-X.
\item
  Let \((\mathbf{W},\mathbf{S})\) be a Coxeter system and \(\mathbf{X},\mathbf{Y}\) two subsets of \(\mathbf{S}\) . Let \(a\in W\) . Show that there exists a unique element \(w\in W_X a W_Y\) of minimal length and that every element
\end{enumerate}

\[
w ^ {\prime} \in W _ {X} a W _ {Y}
\]

can be written in the form \(w' = xwy\) , with \(x \in W_X, y \in W_Y\) and \(l(w') = l(x) + l(w) + l(y)\) (take an element of minimal length in \(W_X aW_Y\) and use Exerc. 1). An element \(w \in W\) is said to be \((X,Y)\) -reduced if it is the element of minimal length in the double coset \(W_X aW_Y\) .

Show that if \(w\) is \((X, \emptyset)\) -reduced then \(l(xw) = l(x) + l(w)\) for all \(x \in W_X\) , and that every element of \(W\) can be written uniquely in the form \(xw\) where \(x \in W_{X}\) and w is \((X, \varnothing)\) -reduced. Show that an element \(w \in W\) is \((X, \varnothing)\) -reduced if and only if \(l(xw) > l(w)\) for all \(x \in X\) (write w in the form \(yw'\) , where \(y \in W_{X}\) and \(w'\) is \((X, \varnothing)\) -reduced).

Show that \(w \in W\) is \((X, Y)\) -reduced if and only if w is both \((X, \varnothing)\) -reduced and \((\varnothing, Y)\) -reduced.

\begin{enumerate}
\def\labelenumi{\arabic{enumi})}
\setcounter{enumi}{3}
\item
  Let \(n\) be an integer \(\geqslant 2\) . For any integer \(i\) with \(1 \leqslant i \leqslant n - 1\) , denote by \(s_i\) the transposition of \(i\) and \(i + 1\) in the set \(\{1, 2, \ldots, n\}\) , and by \(H_i\) the set of \(w \in \mathfrak{S}_n\) such that \(w^{-1}(i) < w^{-1}(i + 1)\) ; put \(S = \{s_1, \ldots, s_{n-1}\}\) . Show that \((\mathfrak{S}_n, S)\) is a Coxeter system and that \(H_i\) is the set of \(w \in \mathfrak{S}_n\) such that \(l(w) < l(s_i w)\) (use Prop. 6 of no. 7).
\item
  Let X be a non-empty set and W a set of permutations of X. Assume given a set R of equivalence relations on X, an element \(x_{0} \in X\) and a map \(\varphi : H \mapsto s_{H}\) from R to W. Denote by \(R_{0}\) the set of \(H \in R\) such that \(s_{\mathrm{H}}(x_{0}) \equiv x_{0} \mod \mathrm{H}'\) for all \(H' \neq H\) in R, and by \(S_{0}\) the set of \(s_{H}\) for H in \(R_{0}\) . We make the following assumptions:
\end{enumerate}

\begin{enumerate}
\def\labelenumi{(\roman{enumi})}
\item
  For any \(\mathbf{H} \in \Re\) , there are two equivalence classes modulo \(\mathbf{H}\) that are permuted by \(s_{\mathrm{H}}\) and \(s_{\mathrm{H}}^{2} = 1\) .
\item
  For all \(\mathbf{H} \in \Re\) and all \(w \in \mathbf{W}\) , the transform \(w(\mathbf{H})\) of \(\mathbf{H}\) by \(w\) is an equivalence relation belonging to \(\Re\) and \(s_{w(\mathbf{H})} = ws_{\mathbf{H}}w^{-1}\) .
\item
  For any \(w \neq 1\) in \(W\) , the set of \(H \in \Re\) such that \(w(x_0) \not\equiv x_0 \mod H\) is finite and meets \(\Re_0\) .
\end{enumerate}

\begin{enumerate}
\def\labelenumi{\alph{enumi})}
\item
  Prove that \((\mathbf{W},\mathbf{S}_0)\) is a Coxeter system (use Prop. 6 of no. 7).
\item
  Prove that the length \(l_{S_0}(w)\) is equal to the number of elements \(H \in \mathfrak{R}\) such that
\end{enumerate}

\[
w (x _ {0}) \not \equiv x _ {0} \bmod . H.
\]

\begin{enumerate}
\def\labelenumi{\alph{enumi})}
\setcounter{enumi}{2}
\tightlist
\item
  Let E be a finite set and X the set of total order relations on X. Denote by W the group of permutations of E, acting in the obvious way on E. Let i and j be distinct elements of E; say that two elements R and R' of X are equivalent mod. \(H_{ij}\) if either \(\mathrm{R}(i,j)\) and \(\mathrm{R}'(i,j)\) or \(\mathrm{R}(j,i)\) and \(\mathrm{R}'(j,i)\) ; and denote the transposition of i and j by \(s_{ij}\) . Let R be the set of equivalence relations of the form \(H_{ij}\) and \(\varphi\) the map from R to W defined by \(\varphi(\mathrm{H}_{ij}) = s_{ij}\) ; finally let \(x_{0}\) be an arbitrary element of X. Show that the objects thus defined satisfy assumptions (i) to (iii), and recover the results of Exerc. 4.
\end{enumerate}

\begin{enumerate}
\def\labelenumi{\arabic{enumi})}
\setcounter{enumi}{5}
\item
  Let E be a set with 6 elements and F the set of structures of the projective line over the field \(F_{5}\) on E. Denote by \(S_{E}\) the group of permutations of E; for any \(\sigma \in S_{E}\) denote by \(\tilde{\sigma}\) the permutation of F induced by \(\sigma\) by transport of structure. Show that there exists a bijection u from E to F, and that the map \(\sigma \mapsto u^{-1}\tilde{\sigma}u\) is an outer automorphism of \(S_{E}\) (if s is a transposition, \(u^{-1}\tilde{s}u\) has three orbits of two elements).
\item
  Construct a group W and two subsets S and S' of W such that (W, S) and (W, S') are isomorphic Coxeter systems, but such that there exists no inner automorphism of W transforming S to S' (use Exerc. 4 and Exerc. 6).
\item
  Construct a group W and two subsets S and S' of W such that (W, S) and (W, S') are non-isomorphic Coxeter systems, one of them being irreducible and the other not (for W take a dihedral group of order 12 generated by \(\{s, s'\}\) where s and \(s'\) are of order 2 and \(s \neq s'\) , and put \(S = \{s, s'\}\) and \(S' = \{(ss')^3, s', s'(ss')^2\}\) ).
\item
  Let (W, S) be a Coxeter system with matrix \((m(s, s'))\) , and let \(W^{+}\) be the subgroup of W consisting of the elements of even length. Let \(s_{0} \in S\) . Put \(g_{s} = ss_{0}\) . Show that the family \((g_{s})_{s \in S - \{s_{0}\}}\) and the relations \(g_{s}^{m(s, s_{0})} = 1\) for \(m(s, s_{0}) \neq \infty\) and \((g_{s} g_{s'}^{-1})^{m(s, s')} = 1\) for \(s, s' \in S - \{s_{0}\}\) and \(m(s, s') \neq \infty\) , form a presentation of \(W^{+}\) . (Let \(H^{+}\) be the group defined by the above presentation. Show that there exists an automorphism \(\alpha\) of \(H^{+}\) , with square the identity, that transforms \(g_{s}\) to \(g_{s}^{-1}\) for all \(s \in S - \{s_{0}\}\) . If \(H_{\alpha}\) is the semi-direct product of \(\{1, -1\}\) and \(H^{+}\) , relative to \(\alpha\) , define mutually inverse homomorphisms \(H_{\alpha} \to W\) and \(W \to H_{\alpha}\) .) Show that if the elements of S are conjugate (cf.~Prop. 3), the group \(W^{+}\) is the commutator subgroup of W (remark that the elements \(g_{s}\) are then commutators).
\item
  Let \(\mathfrak{U}_n\) be the alternating group consisting of the permutations \(w \in \mathfrak{S}_n\) whose signature is equal to \(+1\) . Show that \(\mathfrak{U}_n\) is the commutator subgroup of \(\mathfrak{S}_n\) (use Exerc. 4 and 9). For any integer \(i\) with \(1 \leqslant i \leqslant n - 2\) , put \(u_i = s_i s_{i+1}\) (in the notation of Exerc. 4). Show that the family \((u_i)\) and the relations \(u_1^3 = 1\) , \(u_i^2 = 1\) for \(i \geqslant 2\) , \((u_i u_{i+1})^3 = 1\) for \(1 \leqslant i \leqslant n - 3\) , and \(u_i u_j = u_j u_i\) for \(1 \leqslant i \leqslant n - 4\) and \(i + 2 \leqslant j \leqslant n - 2\) , form a presentation of the group \(\mathfrak{U}_n\) (use Exerc. 9).
\end{enumerate}

*11) Let \((\mathbf{W},\mathbf{S})\) be a Coxeter system. Let \(\Gamma_{\infty}\) be the graph whose set of vertices is \(\mathbf{S}\) , two vertices \(s,s'\) being joined by an edge if and only if \(m(s,s')\neq \infty\) . Let \(\mathrm{S}_{\alpha}\) be the connected components of \(\Gamma_{\infty}\) . Show that \(\mathbf{W}\) can be identified with the free product of the \(\mathrm{W}_{\mathrm{S}_{\alpha}}\) . In particular, every \(w\in \mathbf{W}\) can be written uniquely as a product \(w_{1}\ldots w_{h}\) with \(w_{i}\neq 1\) , \(w_{i}\) belonging to \(\mathrm{W}_{\mathrm{S}_{\alpha_i}}\) , and \(\alpha_{i}\neq \alpha_{i + 1}\) for \(1\leqslant i\leqslant h - 1\) ; show that the length of \(w\) is the sum of the lengths of the \(w_i\).*

\begin{enumerate}
\def\labelenumi{\arabic{enumi})}
\setcounter{enumi}{11}
\tightlist
\item
  Let \((\mathbf{W},\mathbf{S})\) be a Coxeter system such that \(m(s,s')\) is even if \(s\neq s'\) and let \(\mathbf{X}\) be a subset of \(\mathbf{S}\) . Show that there exists a unique homomorphism \(f_{\mathbf{X}}\) from \(\mathbf{W}\) to \(\mathbf{W}_{\mathbf{X}}\) such that \(f_{\mathbf{X}}(s) = s\) for \(s\in \mathbf{X}\) and \(f_{\mathbf{X}}(s) = 1\) for \(s\in \mathbf{S} - \mathbf{X}\) . Deduce that \(\mathbf{W}\) is the semi-direct product of \(\mathbf{W}_{\mathbf{X}}\) and the kernel of \(f_{\mathbf{X}}\) . Show that if \(\mathbf{X}\subset \mathbf{Y}\subset \mathbf{S}\) , there exists a unique homomorphism \(f_{\mathbf{X},\mathbf{Y}}\) from \(\mathbf{W}_{\mathbf{Y}}\) to \(\mathbf{W}_{\mathbf{X}}\) such that \(f_{\mathbf{X}} = F_{\mathbf{X},\mathbf{Y}}\circ f_{\mathbf{Y}}\) and that \(\mathbf{W}\) can be identified with a subgroup of the projective system thus obtained from the \(W_{X}\) as X runs through the filtered set of finite subsets of S.
\end{enumerate}

¶ 13) Let \((W, S)\) be a Coxeter system. For \(s, s' \in S\) , define the sequence \(\mathbf{a}(s, s')\) by means of the following rules:

\begin{enumerate}
\def\labelenumi{(\roman{enumi})}
\item
  if \(ss'\) is of infinite order, \(\mathbf{a}(s, s')\) is the empty sequence;
\item
  if \(ss'\) is of finite order m the sequence \(\mathbf{a}(s, s')\) is of length m and its even (resp. odd) numbered terms are equal to \(s'\) (resp. s).
\end{enumerate}

Denote by \(a(s, s')\) the product of the sequence \(\mathbf{a}(s, s')\) .

\begin{enumerate}
\def\labelenumi{\alph{enumi})}
\item
  Show that the generating set S and the relations \(s^{2}=1\) and \(\mathbf{a}(s,s')=\mathbf{a}(s',s)\) form a presentation of the group W.
\item
  Let \(q\) be an integer \(\geqslant 1\) . Let \(\Sigma_q\) be the set of sequences of \(q\) elements of \(S\) and let \(R_q\) be the smallest equivalence relation on \(\Sigma_q\) for which sequences of the form \((A, \mathbf{a}(s, s'), B)\) and \((A, \mathbf{a}(s', s), B)\) (where \(s, s' \in S\) and \(A\) and \(B\) are sequences of elements of \(S\) ) are equivalent. Let \(\Sigma_q^r\) be the set of sequences \(\mathbf{s} = (s_1, \ldots, s_q)\) such that \(w(\mathbf{s}) = s_1 \ldots s_q\) is of length \(q\) . Show that sequences \(\mathbf{s}, \mathbf{s}' \in \Sigma_q^r\) are equivalent modulo \(R_q\) if and only if \(w(\mathbf{s}) = w(\mathbf{s}')\) (argue by induction on \(q\) and apply Prop. 5 of no. 5).
\item
  Show that a sequence \(\mathbf{s} \in \Sigma_q\) does not belong to \(\Sigma_q^r\) if and only if \(\mathbf{s}\) is equivalent modulo \(R_q\) to a sequence in which two consecutive terms are equal. (Argue by induction on \(q\) and reduce to the case of a sequence \((s_1, \ldots, s_q)\) which does not belong to \(\Sigma_q^r\) but which is such that \((s_1, \ldots, s_{q-1})\) and \((s_2, \ldots, s_q)\) belong to \(\Sigma_{q-1}^r\) ; use Exerc. 1 to show that \(s_1 \ldots s_{q-1} = s_2 \ldots s_q\) and apply b).)
\end{enumerate}

*14) Let \((\mathbf{W},\mathbf{S})\) be a Coxeter system and let \((\Gamma ,f)\) be its Coxeter graph. Let \(k\) be an integer \(\geqslant 3\) . If \(a\) is an edge of \(\Gamma\) , put \(f_{k}(a) = f(a)\) if \(f(a)\neq \infty\) and \(f_{k}(a) = k\) if \(f(a) = \infty\) . Let \((\mathrm{W}_k,\mathrm{S})\) be a Coxeter system whose Coxeter graph is equal to \((\Gamma ,f_k)\) (Chap. V, §4, no. 3, Cor. of Prop. 4). Show that there exists a unique homomorphism \(\varphi_{k}\) from W to \(\mathbf{W}_k\) inducing the identity on S. Show that if \(k\) divides \(k^{\prime}\) , there exists a unique homomorphism \(\varphi_{k,k^{\prime}}\) from \(\mathbf{W}_{k^{\prime}}\) to \(\mathbf{W}_k\) such that \(\varphi_{k} = \varphi_{k,k^{\prime}}\circ \varphi_{k^{\prime}}\) . Show that the homomorphism \((\varphi_{k})\) from W to the projective limit of the \(\mathbf{W}_k\) is injective (use Exerc. 13) \(c)\) , but in general not surjective (for example, in the case of the infinite dihedral group).*

\begin{enumerate}
\def\labelenumi{\arabic{enumi})}
\setcounter{enumi}{14}
\tightlist
\item
  Let A be a set and C a subset of \(\mathfrak{P}(A)\) . The elements of C are called the chambers of A and a subset F of a chamber C is called a facet, the codimension of F in C being the cardinal of C-F. A facet F is said to be a panel of C if F is of codimension 1 in C. Two chambers C and C' are said to be adjoining if they have a common panel F: then either \(C = C'\) or \(F = C \cap C'\) . A gallery of length n is a sequence \(\Gamma = (C_0, C_1, \ldots, C_n)\) of \(n + 1\) chambers such that \(C_i\) and \(C_{i+1}\) are adjoining for \(0 \leqslant i \leqslant n - 1\) . Then \(C_0\) and \(C_n\) are called the ends of \(\Gamma\) . The gallery \(\Gamma\) is called injective if \(C_i \neq C_{i+1}\) for \(0 \leqslant i \leqslant n - 1\) and is called minimal if there is no gallery with the same ends and length \textless{} n.
\end{enumerate}

The set A (together with C) is a building if every element of A belongs to at least one chamber and if any two chambers are the ends of a gallery. The distance between two chambers C and \(C'\) is the length \(d(\mathrm{C},\mathrm{C}')\) of a minimal gallery with ends C and \(C'\) .

A sub-building of a building A is a subset D of A such that D together with \(\mathfrak{C}\cap\mathfrak{P}(D)\) is a building.

\begin{enumerate}
\def\labelenumi{\alph{enumi})}
\item
  Show that if A is a building, a facet has the same codimension in every chamber containing it; this allows us to speak of the codimension of a facet and of a panel of A without reference to a particular chamber. A morphism of a building B to A is a map f from B to A such that the restriction of f to every chamber C of B is a bijection from C to a chamber \(f(\mathrm{C})\) of A. Show that the image under f of a facet of B is a facet of the same codimension.
\item
  A building is called an apartment if every panel is contained in exactly two chambers. Show that if A is an apartment, every automorphism of A (i.e.~every permutation of A preserving C) that leaves fixed all the points of a chamber is the identity. More generally, let \(\varphi\) be an endomorphism of A and let C be a chamber of A such that \(\varphi(a)=a\) for all \(a\in C\) . Let \((\mathrm{C},\mathrm{C}_{1},\ldots,\mathrm{C}_{n})\) be a gallery of A. Show that, either the gallery \((\mathrm{C},\varphi(\mathrm{C}_{1}),\ldots,\varphi(\mathrm{C}_{n}))\) is not injective, or \(\varphi(a)=a\) for all a belonging to the union of the \(C_{i}\) .
\end{enumerate}

\begin{enumerate}
\def\labelenumi{\arabic{enumi})}
\setcounter{enumi}{15}
\tightlist
\item
  Let \((W,S)\) be a Coxeter system. For \(s \in S\) , denote by \(W^{(s)}\) the subgroup \(W_{S-\{s\}}\) of W, by A the set of subsets of W of the form \(wW^{(s)}\) for \(w \in W\) and \(s \in S\) , and by C the set of subsets of A of the form \(C_w = \{wW^{(s)} \mid s \in S\}\) for \(w \in W\) , which we call the chambers of A (Exerc. 15).
\end{enumerate}

\begin{enumerate}
\def\labelenumi{\alph{enumi})}
\item
  Show that the map \(w \mapsto \mathbf{C}_w\) is a bijection from W onto \(\mathfrak{C}\) .
\item
  Show that two distinct chambers \(C_{w}\) and \(C_{w'}\) are adjoining if and only if there exists \(s \in S\) such that \(w' = ws\) . Deduce that A (together with C) is an apartment (Exerc. 15), which we call the apartment of (W, S). Show that the length of a minimal gallery with ends \(C_{w}\) and \(C_{w'}\) is equal to \(l_{S}(w^{-1}w')\) .
\item
  Let \(\mathfrak{F}\) be the set of facets of A and let \(F \in \mathfrak{F}\) . Show that there exist a unique subset X of S and an element \(w \in W\) such that \(wW_X = \bigcap_{i \in F} i\) . Then F is said to be of type X. Show that the codimension of F is equal to the cardinal of X. Show that the map \(j: F \mapsto \bigcap_{i \in F} i\) is a strictly decreasing bijection (with respect to inclusion) from \(\mathfrak{F}\) to the set of subsets of W of the form \(wW_X\) for \(w \in W\) and \(X \subset S\) . Show that every facet of type X contains a unique facet of type Y for every Y such that \(X \subset Y \subset S\) .
\item
  Let W act on A be left translations and put \(C = C_{e}\) . Show that W acts simply-transitively \(^{8}\) on C. Let \(C_{1}, \ldots, C_{n}\) be chambers of A. Show that the following conditions are equivalent:
\end{enumerate}

\begin{enumerate}
\def\labelenumi{(\roman{enumi})}
\item
  the sequence \(\Gamma = (\mathrm{C}_0 = \mathrm{C},\mathrm{C}_1,\dots ,\mathrm{C}_n)\) is an injective gallery;
\item
  there exists a sequence \(\mathbf{s} = (s_1, \ldots, s_n)\) of elements of S such that \(\mathrm{C}_j = t_j(\mathrm{C}_{j-1})\) for \(1 \leqslant j \leqslant n\) , where \(t_j\) is the element of \(\Phi(\mathbf{s})\) defined in no. 4, formula (11).
\end{enumerate}

Show that if these conditions are satisfied, the sequence-s is unique and \(C_{n} = s_{1} \ldots s_{n}(C)\) . Show that the gallery \(\Gamma\) is minimal if and only if the sequence \(\mathbf{s}(\Gamma) = \mathbf{s}\) is a reduced decomposition of \(w = s_{1} \ldots s_{n}\) .

\begin{enumerate}
\def\labelenumi{\alph{enumi})}
\setcounter{enumi}{4}
\tightlist
\item
  Let T be the union of the conjugates of S. For \(t \in T\) , the set \(L_{t}\) of points of A invariant under t is called the wall defined by t. Show that \(L_{t}\) is a union of panels and that a panel F is contained in \(L_{t}\) if and only if \(j(\mathrm{F})\) is of the form \(wW_{\{s\}}\) with \(t = wsw^{-1}\) . Deduce that, for any panel F, there is a unique element \(t = t(\mathrm{F}) \in \mathrm{T}\) such that \(F \subset L_{t}: L_{t}\) is called the support of F.
\end{enumerate}

Show that if \(w(\mathrm{L}_t) = \mathrm{L}_t\) (for \(w \in \mathrm{W}\) ), then \(w = 1\) or \(w = t\) .

\(f)\) Let \(w', w'' \in W\) . Put \(C' = w'(C), C'' = w''(C)\) and let \(\Gamma = (C_0 = C', C_1, \ldots, C_n = C'')\) be an injective gallery with ends \(C'\) and \(C''\) . Let \(t_j\) be the element of \(T\) defining the support wall of the panel \(C_j \cap C_{j-1}\) (for \(1 \leqslant j \leqslant n\) ). Show that the sequence \(\Psi(T) = (w'^{-1} t_j w')_{1 \leqslant j \leqslant n}\) coincides with the sequence \(\Phi(s(w'^{-1}(\Gamma))\) . For \(t \in T\) , let \(n(\Gamma, t)\) be the number of times \(w'^{-1} tw'\) appears in \(\Psi(\Gamma)\) . Deduce from Lemma 1 of no. 4 that the number \((-1)^{n(\Gamma, t)}\) depends only on \(C'\) and \(C''\) and not on \(\Gamma\) : we denote it by \(\eta(C', C'', t)\) . Show that the relation \(\eta(C', C'', t) = 1\) is an equivalence relation between \(C'\) and \(C''\) that has two equivalence classes that are permuted by \(t\) . Denote by \(\mathfrak{C}^+(t)\) the equivalence class containing \(C\) and by \(\mathfrak{C}^-(t)\) the other.

Show that, for \(w \in \mathbf{W}\) and \(s \in S\) , the chamber \(w(\mathbf{C})\) belongs to \(\mathfrak{C}^+(s)\) if and only if \(l(sw) > l(w)\) .

\(g)\) Let \(\mathrm{A}^{+}(t)\) (resp. \(\mathrm{A}^{-}(t)\) ) be the union of the chambers belonging to \(\mathfrak{C}^{+}(t)\) (resp. \(\mathfrak{C}^{-}(t)\) ) (for \(t \in \mathrm{T}\) ). Show that \(\mathrm{A}^{+}(t) \cap \mathrm{A}^{-}(t) = \mathrm{L}_{t}\) . (To show that \(\mathrm{A}^{+}(t) \cap \mathrm{A}^{-}(t) \subset \mathrm{L}_{t}\) , reduce to the case \(t \in \mathrm{S}\) . If \(a \in \mathrm{A}^{+}(t) \cap \mathrm{A}^{-}(t)\) , put \(a = w\mathrm{W}^{(s)}\) with \(s \in \mathrm{S}\) and \(w\) being \((\varnothing, \mathrm{S} - \{s\})\) -reduced (Exerc. 3). If \(w(\mathrm{C}) \in \mathfrak{C}^{-}(t)\) , then \(l(tw) < l(w)\) and \(w = ts_{1}, \ldots, s_{q}\) with \(s_{j} \in \mathrm{S}\) and \(l(w) = q + 1\) . Since \(a \in \mathrm{A}^{-}(t)\) , there exists \(x \in \mathrm{W}^{(s)}\) such that \(wx(\mathrm{C}) \in \mathfrak{C}^{+}(t)\) . Then

\[
l (t w x) = 1 + l (w x) = 1 + l (w) + l (x),
\]

but since \(twx = s_1 \ldots s_q x\) this is a contradiction. Thus, \(w(\mathbf{C}) \in \mathfrak{C}^+(t)\) and \(l(tw) = 1 + l(w)\) . If \(x \in \mathrm{W}^{(s)}\) is such that \(l(twx) < l(wx)\) , deduce from Exerc. 1 that \(twx = wx'\) with \(x' \in \mathrm{W}^{(s)}\) , and hence that \(ta = a\) and \(a \in \mathrm{L}_t\) .

The subsets \(\mathbf{A}^{+}(t)\) and \(\mathbf{A}^{-}(t)\) are called the halves of A determined by the wall \(L_{t}\) . Two points of A are said to be on the same side (resp. strictly on opposite sides) of \(L_{t}\) if they belong (resp. do not both belong) to one of the two halves. Any facet is contained in one of the halves determined by \(L_{t}\) . If two facets are contained in different halves, they are said to be on opposite sides of \(L_{t}\) , or to be separated by \(L_{t}\) .

\(h\) ) Let \(w \in \mathbf{W}\) . Show that \(l(w)\) is equal to the number of walls separating \(\mathbf{C}\) and \(w(\mathbf{C})\) .

\begin{enumerate}
\def\labelenumi{\roman{enumi})}
\tightlist
\item
  Show that the map \(\varphi\) that takes the half \(A^{+}(t)\) (resp. \(A^{-}(t)\) ) to \((1,t)\) (resp. \((-1,t)\) ) is a bijection from the set \(\mathfrak{M}\) of halves of \(A\) onto the set \(R = \{1, - 1\} \times T\) (cf.~no. 4). With the notation of Lemma 1 of no. 4, show that \(\varphi(w(M)) = U_w(\varphi(M))\) for all \(w\in W\) and \(M\in \mathfrak{M}\) .
\end{enumerate}

\begin{enumerate}
\def\labelenumi{\arabic{enumi})}
\setcounter{enumi}{16}
\item
  We retain the notation of Exerc. 16 and assume that W is finite. Let H be the set of walls of A. To any \(H \in H\) , we associate the half \(H^{+}\) of A determined by H containing C. Show that the elements of H can be numbered in such a way that the map \(j \mapsto \bigcap_{i \leqslant j} H_{i}^{+}\) is strictly decreasing on the interval {[}1, \(\operatorname{Card}(\mathfrak{H})\) {]}. (Consider the family \(\mathfrak{F}\) of intersections of the sets \(\mathrm{H}^{+}\) ordered by inclusion and consider a strictly decreasing sequence \((\mathrm{F}_{0},\ldots,\mathrm{F}_{q})\) of elements of \(\mathfrak{F}\) of maximal length. For any \(\mathrm{H}\in\mathfrak{H}\) , there exists an \(i\) such that \(\mathrm{H}^{+}\supset\mathrm{F}_{j}\) for \(j\geqslant i\) and \(\mathrm{H}^{+}\not\supset\mathrm{F}_{i-1}\) : show that \(\mathrm{F}_{i}=\mathrm{H}^{+}\cap\mathrm{F}_{i-1}\) .)
\item
  Let A be an apartment (Exerc. 15). A folding of A is an endomorphism \(\pi\) of A such that \(\pi^{2} = \pi\) and such that every chamber of A is the image under \(\pi\) of 0 or 2 chambers.
\end{enumerate}

\begin{enumerate}
\def\labelenumi{\alph{enumi})}
\item
  Let \(\pi\) be a folding of A and C a chamber of A such that \(\pi(C) = C\) . For any neighbouring chamber \(C'\) of C, either \(\pi(C') = C'\) or \(\pi(C') = C\) ; if \(a \in C\) we have \(\pi(a) = a\) . Show that, if \((C_0 = C, C_1, \ldots, C_n)\) is a gallery, either \(\pi(C_i) = C_i\) for all \(i\) or \((C_0, \pi(C_1), \ldots, \pi(C_n))\) is not minimal (and in fact has two consecutive equal chambers). Deduce that any minimal gallery whose extremities are invariant under \(\pi\) is invariant under \(\pi\) . If \((C = C_0, C_1, \ldots, C_n)\) is a minimal gallery and if \(\pi(C_n) \neq C_n\) , there exists an index \(i\) with \(0 \leqslant i < n\) such that \(\pi(C_j) = C_j\) for \(0 \leqslant j \leqslant i\) and \(\pi(C_j) \neq C_j\) for \(i < j \leqslant n\) .
\item
  Let \(C_{1}\) and \(C_{2}\) be two distinct neighbouring chambers and \(\pi, \pi'\) two foldings of A. Assume that \(\pi(C_{0}) = C_{1}\) and \(\pi'(C_{1}) = C_{0}\) . Let C be a chamber, and consider the following three conditions:
\end{enumerate}

\[
\pi (\mathrm{C}) = \mathrm{C};\tag{1}
\]

\[
d (\mathrm{C}, \mathrm{C} _ {1}) <   d (\mathrm{C}, \mathrm{C} _ {2});\tag{2}
\]

\[
\pi^ {\prime} (\mathrm{C}) \neq \mathrm{C}.\tag{3}
\]

Show that (1) \(\Longrightarrow\) (2) \(\Longrightarrow\) (3) and deduce that the three conditions are equivalent. Show that \(\pi\) (resp. \(\pi'\) ) is the unique folding of A taking \(C_2\) (resp.

\(C_{1}\) to \(C_{1}\) (resp. \(C_{2}\) ) (assume that (2) is satisfied and let \((C_{1}, C_{2}', \ldots, C_{n}' = C)\) be a minimal gallery: show that \(\pi'(C_{j+1}')\) is the unique chamber distinct from \(\pi'(C_{j}')\) and containing the panel \(\pi'(C_{j}' \cap C_{j+1}')\) . Show that \(\pi(\mathfrak{C})\) and \(\pi'(\mathfrak{T})\) form a partition of the set T of chambers of A and that \(\pi(a) = \pi'(a) = a\) for all \(a \in \pi(A) \cap \pi'(A)\) . Show that the map from A to itself that coincides with \(\pi'\) on \(\pi(A)\) and with \(\pi\) on \(\pi'(A)\) is an involutive automorphism of A. It is called the reflection with respect to the panel \(C_{1} \cap C_{2}\) . Show that it is the only non-trivial automorphism of A leaving fixed all the points of \(C_{1} \cap C_{2}\) (use Exerc. 15 b)).

\begin{enumerate}
\def\labelenumi{\alph{enumi})}
\setcounter{enumi}{2}
\tightlist
\item
  Suppose that A is the apartment associated to a Coxeter system (W,S) and retain the notation of Exerc. 16. Let \(C_{1}\) and \(C_{2}\) be two neighbouring chambers and let t be the element of T such that the wall \(L_{t}\) is the support of \(C_{1} \cap C_{2}\) . Let \(M_{j}\) be the half of A determined by \(L_{t}\) and containing \(C_{j}\) (for j = 1,2). Show that the map \(\pi\) defined by \(\pi(a) = a\) if \(a \in M_{1}\) and \(\pi(a) = t(a)\) if \(a \in M_{2}\) is a folding of A such that \(\pi(C_{2}) = C_{1}\) and that the reflection with respect to the panel \(C_{1} \cap C_{2}\) is the map \(a \mapsto t(a)\) .
\end{enumerate}

\begin{enumerate}
\def\labelenumi{\arabic{enumi})}
\setcounter{enumi}{18}
\tightlist
\item
  Let A be an apartment. Assume that, for any two distinct neighbouring chambers \(C_{1}\) and \(C_{2}\) , there is a folding (Exerc. 18) of A taking \(C_{1}\) to \(C_{2}\) . Let C be a chamber of A and \((\mathrm{C}_{i})_{i\in I}\) the family of chambers neighbouring C and distinct from C. Denote by \(s_{i}\) the reflection with respect to the panel \(C\cap C_{i}\) (Exerc. 18 b)). Put \(S=\{s_{i}\mid i\in I\}\) and denote by W the group of automorphisms of A generated by the \(s_{i}\) .
\end{enumerate}

\begin{enumerate}
\def\labelenumi{\alph{enumi})}
\item
  Show that, for any chamber \(C'\) , there exists \(w \in W\) such that \(C' = w(C)\) (argue by induction on \(d(C, C')\) ).
\item
  Show that (W,S) is a Coxeter system (for \(i\in I\) ,put
\end{enumerate}

\[
\mathrm{P} _ {s _ {i}} = \left\{w \in \mathrm{W} \mid w (\mathrm{C}) \subset \pi_ {i} (\mathrm{A}) \right\},
\]

where \(\pi_{i}\) is the folding taking \(C_{i}\) to C, and show that the assumptions of Prop. 6 are satisfied: to prove condition (C'), remark that if \(w \in P_{s_{i}}\) and \(ws_{j} \notin P_{s_{i}}\) , then

\[
w (\mathrm{C}) \cap w s _ {j} (\mathrm{C}) \subset \pi_ {i} (\mathrm{A}) \cap s _ {i} \pi_ {i} (\mathrm{A}).
\]

Since \(w(\mathbf{C})\) and \(ws_{j}(\mathbf{C})\) are neighbouring, it follows that \(s_i = ws_jw^{-1}\) (Exerc. 18 b)).

\begin{enumerate}
\def\labelenumi{\alph{enumi})}
\setcounter{enumi}{2}
\item
  Let F be a facet of the chamber C. Show that the stabilizer \(W_{F}\) of F in W is generated by the \(s_{i} \in S\) such that \(F \subset C \cap C_{i}\) (let \(w \in W_{F}\) with \(l_{S}(w) > 1\) and let \(i \in I\) be such that \(w = s_{i}w'\) with \(l(w') = l(w) - 1\) ; by Prop. 6, \(w' \in P_{s_{i}}\) , hence \(w(\mathrm{C}) \subset s_{i}\pi_{i}(\mathrm{A})\) , \(F \subset \pi_{i}(\mathrm{A}) \cap s_{i}\pi_{i}(\mathrm{A})\) and \(s_{i} \in W_{F}\) ). In particular, \(w(\mathrm{C}) = \mathrm{C}\) if and only if w = 1.
\item
  Show that the map \(a \mapsto W_{\{a\}}\) is an isomorphism from the apartment A to the apartment associated to (W, S) (Exerc. 16), compatible with the action of W.
\end{enumerate}

\begin{enumerate}
\def\labelenumi{\arabic{enumi})}
\setcounter{enumi}{19}
\tightlist
\item
  Let A be a building and S a set. Then A is said to be numbered by S if one is given a map f from A to S such that, for any chamber C of A, the restriction of f to C is a bijection from C to S. If F is a facet of A, \(f(F)\) is called the type of F. Let A be a numbered building. An endomorphism \(\varphi\) of A is called allowed if a and \(\varphi(a)\) are of the same type for all \(a \in A\) .
\end{enumerate}

\begin{enumerate}
\def\labelenumi{\alph{enumi})}
\item
  Let \(\varphi\) be an endomorphism of A. Show that, if there exists a chamber C of A such that a and \(\varphi(a)\) are of the same type for all \(a \in C\) , then \(\varphi\) is allowed. Show that if A is an apartment and \(\pi\) is a folding of A (Exerc. 18), then \(\pi\) is an allowed endomorphism.
\item
  A subset D of A is called convex if, for all \(a \in A - D\) , there exists an allowed endomorphism \(\varphi\) of A such that \(\varphi(x) = x\) for all \(x \in D\) and \(\varphi(a) \neq a\) . Show that any intersection of convex sets is convex and that, for any subset D of A, there exists a smallest convex subset containing D: this is called the convex hull of D and denoted by \(\Gamma(D)\) .
\end{enumerate}

\begin{enumerate}
\def\labelenumi{\arabic{enumi})}
\setcounter{enumi}{20}
\tightlist
\item
  Let \((W,S)\) be a Coxeter system and A the associated apartment (cf.~Exerc. 16, of which we retain the notation).
\end{enumerate}

\begin{enumerate}
\def\labelenumi{\alph{enumi})}
\item
  Show that there exists a unique numbering of A (called the canonical numbering) for which the type of a facet F is that defined in Exerc. 16 c). We shall always consider A to be equipped with this numbering.
\item
  Show that the allowed automorphisms of \(\mathbf{A}\) are the operations of \(\mathbf{W}\) .
\item
  Let D be a subset of A containing at least one chamber. Show that the following conditions are equivalent:
\end{enumerate}

\begin{enumerate}
\def\labelenumi{(\roman{enumi})}
\item
  D is the intersection of the halves of A (Exerc. 16 g)) that contain D;
\item
  D is convex;
\item
  whenever two facets \(F_{1}\) and \(F_{2}\) are contained in D, the convex hull of \(F_{1} \cup F_{2}\) is contained in D;
\item
  whenever a chamber \(\mathbf{C}_1\) and a facet \(\mathbf{F}\) are contained in \(\mathrm{D}\) and \((C_1, \ldots, C_n)\) is a gallery of minimal length such that \(\mathbf{F} \subset \mathbf{C}_n\) , we have \(C_i \subset D\) for \(1 \leqslant i \leqslant n\) .
\end{enumerate}

(To show that (iii) \(\Longrightarrow\) (iv), use Exerc. 15 \(b\) ). To show that (iv) \(\Longrightarrow\) (i), argue by contradiction. Let \(D'\) be the intersection of the halves of A containing D; let \(a \in D' - D\) , let \(C_0\) be a chamber contained in D and let \((C_0, C_1, \ldots, C_n)\) be a gallery of minimal length such that \(a \in C_n\) . Then \(C_j \subset D'\) for all \(j\) . Show that there exists an integer \(j\) with \(0 \leqslant j < n\) such that \(C_j \subset D\) and \(C_{j+1} \not\subset D\) . Let M (resp. \(M'\) ) be the half of A determined by the support wall of the panel \(C_j \cap C_{j+1}\) and containing \(C_j\) (resp. \(C_{j+1}\) ). Show that D \(\not\subset M\) . Let \(b \in D \cap (A - M)\) and let \(\Gamma = (C_j, C_1', \ldots, C_p')\) be a gallery of smallest possible length such that \(b \in C_p'\) . Then \(C_k' \subset D\) for \(1 \leqslant k \leqslant p\) and \(C_p' \subset M'\) . Let \(\pi\) be the folding of A with image \(M'\) (Exerc. 18 \(c\) ): then \(\pi(C_j) = C_{j+1}\) and the gallery \(\pi(\Gamma)\) is not injective (Exerc. 18 \(a\) ). If \(\Gamma' = (C_{j+1}, C_2', \ldots, C_{p-2}', C_p')\)

is the gallery obtained from \(\pi(\Gamma)\) by suppressing one of the two equal consecutive chambers, the gallery \((\mathrm{C}_{j}, \mathrm{C}_{j+1}, \mathrm{C}_{2}', \ldots, \mathrm{C}_{p-2}', \mathrm{C}_{p}')\) is minimal by the definition of \(\Gamma\) . Deduce from (iv) that \(\mathrm{C}_{j+1} \subset \mathrm{D}\) . Contradiction.)

\begin{enumerate}
\def\labelenumi{\arabic{enumi})}
\setcounter{enumi}{21}
\tightlist
\item
  We retain the notation of Exerc. 16 and 21.
\end{enumerate}

\begin{enumerate}
\def\labelenumi{\alph{enumi})}
\item
  Let \(t \in \mathbf{T}\) and \(w \in \mathbf{W}\) . Show that the chambers \(\mathbf{C}\) and \(w(\mathbf{C})\) are separated by the wall \(\mathbf{L}_t\) if and only if \(l(tw) < l(w)\) (use the folding determined by the half \(\mathbf{A}^+(t)\) ).
\item
  Let \(w_0 \in W\) . Show that the following conditions are equivalent:
\end{enumerate}

\[
l (w w _ {0}) = l (w _ {0}) - l (w) \quad \text {   for   all   } w \in \mathrm{W};\tag{i}
\]

\[
l (t w _ {0}) <   l (w _ {0}) \quad \text {   for   all   } t \in \mathrm{T};\tag{ii}
\]

whenever \(t \in \mathrm{T}\) , the chambers \(\mathrm{C}\) and \(w_0(\mathrm{C})\) are separated by the wall \(\mathbf{L}_t\) (iii)

(Use Exerc. 16 h) to show that (iii) \(\Longrightarrow\) (i).)

Show that such an element \(w_{0}\) is unique and exists if and only if W is finite. It is then the element of greatest length of W and is characterized by

\[
l (s w _ {0}) <   l (w _ {0}) \quad \text {   for   all   } s \in S.\tag{iv}
\]

Moreover, \(w_0^2 = 1\) , \(w_0Sw_0 = S\) and \(l(w_0) = \mathrm{Card}(T)\) .

\begin{enumerate}
\def\labelenumi{(\alph{enumi})}
\setcounter{enumi}{2}
\tightlist
\item
  Assume that W is finite. Show that, for any chamber \(C_{0}\) , there exists a unique chamber \(-C_{0}\) such that \(C_{0} \cup (-C_{0})\) is not contained in any half of A. Show that there exists a unique involutive automorphism \(\varphi\) of A (not necessarily allowed such that \(\varphi(C_{0}) = -C_{0}\) for any chamber \(C_{0}\) and that \(\varphi(L) = L\) for any wall L of A. We put \(\varphi(a) = -a\) for \(a \in A\) . If F is a facet, the facet \(-F = \varphi(F)\) will be said to be opposed to F.
\end{enumerate}

\begin{enumerate}
\def\labelenumi{\alph{enumi})}
\setcounter{enumi}{3}
\tightlist
\item
  Let \(C_{0}\) be a chamber of A and F a partition of \(C_{0}\) . Show that the convex hull of \(C_{0} \cup (-F)\) is the half of A determined by the wall L that is the support of F and that contains \(C_{0}\) .
\end{enumerate}

\begin{enumerate}
\def\labelenumi{\arabic{enumi})}
\setcounter{enumi}{22}
\item
  We retain the notation of Exerc. 16. Let \(\operatorname{Aut}(A)\) be the group of automorphisms of the apartment \(A\) . Show that if \(\varphi \in \operatorname{Aut}(A)\) , then \(\varphi\) permutes the walls of \(A\) , and \(\varphi t\varphi^{-1} \in T\) for all \(t \in T\) (use Exerc. 18). Deduce that \(W\) (identified with a subgroup of \(\operatorname{Aut}(A)\) ) is normal and that \(\operatorname{Aut}(A)\) is the semi-direct product of the subgroup \(E\) of automorphisms preserving the chamber \(C\) by \(W\) . Show that the action of \(\operatorname{Aut}(A)\) on \(W\) defines an isomorphism from \(E\) to the group of automorphisms of the Coxeter system \((W, S)\) , or to that of the Coxeter graph of \((W, S)\) (cf.~also Exerc. 19).
\item
  A structured building is a building I equipped with a set U of sub-buildings satisfying the following conditions:
\end{enumerate}

(SB 1) The sub-buildings \(A \in \mathfrak{U}\) are apartments.

(SB 2) Any two chambers of I are contained in at least one element of U. (SB 3) Whenever \(A_{1}, A_{2} \in U\) are such that \(A_{1} \cap A_{2}\) contains a chamber, there exists an isomorphism from \(A_{1}\) to \(A_{2}\) leaving fixed the points of \(A_{1} \cap A_{2}\) .

Let \((\mathbf{I},\mathfrak{U})\) be a structured building. The elements of \(\mathfrak{U}\) are called the apartments of \((\mathbf{I},\mathfrak{U})\) , or simply those of \(\mathbf{I}\) .

\begin{enumerate}
\def\labelenumi{\alph{enumi})}
\item
  Show that any two apartments of I are isomorphic. Let C be a chamber of I and A an apartment of I containing C. Show that there exists a unique endomorphism \(\rho\) of I (called the retraction with centre C of I onto A) such that \(\rho(a)=a\) for all \(a\in A\) and that, for any apartment \(A'\) containing C, the restriction of \(\rho\) to \(A'\) is an isomorphism from \(A'\) to A (remark that, by (SB 2) and Exerc. 15 b), for any apartment \(A'\) containing C, there exists a unique isomorphism \(\rho_{A'}\) of \(A'\) leaving all the points of C fixed). Show that \(\rho^{2}=\rho\) and that \(\rho^{-1}(C)=C\) .
\item
  Let A be an apartment of I, \(C_{0}\) a chamber and F a facet contained in A. Let \((C_{0}, C_{1}, \ldots, C_{n})\) be a gallery of smallest possible length such that \(F \subset C_{n}\) . Show that \(C_{i} \subset A\) for \(1 \leqslant i \leqslant n\) (argue by contradiction: if \(C_{i} \subset A\) and \(C_{i+1} \notin A\) , consider the retraction of I onto A with centre the chamber of A distinct from \(C_{i}\) and containing \(C_{i} \cap C_{i+1}\) ).
\item
  Let A be an apartment of I, C a chamber of A, F a panel of C, C' a chamber of I and \(\Gamma = (C_0, \ldots, C_n = C')\) a gallery of smallest possible length such that \(F \subset C_0\) . Show that the retraction \(\rho\) of I onto A with centre C transforms \(\Gamma\) into a gallery
\end{enumerate}

\[
\left(\mathrm{C} _ {0} ^ {\prime}, \dots , \mathrm{C} _ {n} ^ {\prime} = \rho (\mathrm{C} ^ {\prime})\right)
\]

of smallest possible length such that \(F \subset C_{0}^{\prime}\) (consider an apartment \(A^{\prime}\) containing \(C^{\prime}\) and F and apply b) and the fact that the restriction of \(\rho\) to \(A^{\prime}\) is an isomorphism).

\begin{enumerate}
\def\labelenumi{\alph{enumi})}
\setcounter{enumi}{3}
\tightlist
\item
  Let A be an apartment, \(C_{1}\) and \(C_{2}\) two distinct neighbouring chambers contained in A, \(C'\) a chamber containing \(C_{1} \cap C_{2}\) and distinct from \(C_{1}\) and \(C_{2}\) , and \(A_{i}'\) an apartment containing \(C_{i}\) and \(C'\) (for i = 1, 2). Let \(\varphi_{i}\) (resp. \(\psi_{i}\) ) be the retraction of I onto A (resp. \(A_{i}'\) ) with centre \(C_{i}\) (resp. \(C'\) ) and let \(\rho_{i}\) be the restriction of \(\varphi_{i} \circ \psi_{i}\) to A. Let C be a chamber of A; show that if
\end{enumerate}

\[
d (\mathrm{C}, \mathrm{C} _ {1}) \leqslant d (\mathrm{C}, \mathrm{C} _ {2}),
\]

then \(\rho_{1}(a) = a\) for all \(a \in \mathbb{C}\) and \(d(\mathbb{C}, \mathbb{C}_{1}) < d(\mathbb{C}, \mathbb{C}_{2})\) (consider a minimal gallery \(\Gamma\) with extremities \(\mathbb{C}\) and \(\mathbb{C}_{1}\) and apply \(c\) ) to show that \(\rho_{1}(\Gamma)\) is minimal; then use Exerc. 15 b)). Show that, if \(d(\mathbb{C}, \mathbb{C}_{2}) < d(\mathbb{C}, \mathbb{C}_{1})\) , then \(\rho_{1}(\mathbb{C}) \neq \mathbb{C}\) and \(\rho_{1}^{2}(\mathbb{C}) = \rho_{1}(\mathbb{C})\) (take a minimal gallery \(\Gamma\) with extremities \(\mathbb{C}\) and \(\mathbb{C}_{2}\) and use \(c\) ) to show that \(\rho_{1}(\Gamma)\) is a gallery of smallest possible length having one extremity equal to \(\rho_{1}(\mathbb{C})\) and the other containing \(\mathbb{C}_{1} \cap \mathbb{C}_{2}\) ; deduce that \(d(\mathbb{C}_{1}, \rho_{1}(\mathbb{C}_{1})) \leqslant d(\mathbb{C}_{2}, \rho_{1}(\mathbb{C}_{1}))\) .

Show that \(\rho_{i}\) is a folding (Exerc. 19) of A (show that \(A = \rho_{1}(A) \cup \rho_{2}(A)\) and define an involutive automorphism \(\sigma\) of A by putting \(\sigma(a) = \rho_{2}(a)\) if \(a \in \rho_{1}(A)\) and \(\sigma(a) = \rho_{1}(a)\) if \(a \in \rho_{2}(A)\) ; show that if C is a chamber contained in \(\rho_{1}(A)\) , then \(\rho_{1}^{-1}(C) = \{C, \rho_{2}(C)\}\) ).

\begin{enumerate}
\def\labelenumi{\alph{enumi})}
\setcounter{enumi}{4}
\tightlist
\item
  Let I be a spacious structured building, that is to say that every panel is contained in at least three chambers. Show that there exists a Coxeter system (W,S), unique up to isomorphism, such that the apartments of I are isomorphic to the apartment \(A_{0}\) associated to (W,S) (use d) and Exerc. 19). Let A be an apartment of I and \(\varphi\) an isomorphism from \(A_{0}\) to A. Show that there exists a unique numbering of I, with values in S, such that the types of a and \(\varphi(a)\) are equal for all \(a \in A_{0}\) (choose a chamber C of A and show, by using b), that, if \(A'\) and \(A''\) are two apartments of I containing C, the numberings of \(A'\) and \(A''\) extending that determined on C coincide on \(A' \cap A''\) ).
\end{enumerate}

We shall say that the Coxeter system (W,S) and the numbering thus obtained are adapted to I.

Show that the retractions introduced in a) are allowed endomorphisms. Show that a subset of I contained in an apartment A of I is convex in I if and only if it is convex in A (Exerc. 20).

\(f\) ) We retain the notation of \(e\) ) and let \(\mathfrak{A}\) be the set of allowed isomorphisms from \(\mathbf{A}_0\) to the various apartments of I. Show that, if \(\varphi, \psi \in \mathfrak{A}\) and if C is a chamber and F a facet of I contained in \(\varphi(\mathbf{A}_0) \cap \psi(\mathbf{A}_0)\) , there exists an element \(w \in W\) such that \(\psi^{-1}(C) = w\varphi^{-1}(C)\) and \(\psi^{-1}(F) = w\varphi^{-1}(F)\) (consider an isomorphism \(\lambda\) from \(\varphi(\mathbf{A}_0)\) to \(\psi(\mathbf{A}_0)\) leaving fixed the points of \(\varphi(\mathbf{A}_0) \cap \psi(\mathbf{A}_0)\) and apply Exerc. 21 \(b\) ) to the automorphism \(\psi^{-1} \circ \lambda \circ \varphi^{-1}\) of \(\mathbf{A}_0\) ).

\begin{enumerate}
\def\labelenumi{\arabic{enumi})}
\setcounter{enumi}{24}
\tightlist
\item
  Let I be a set and let \(\mathfrak{F}\) be the set of finite subsets of I. For any \(A \in \mathfrak{F}\) , put \(\varepsilon(A) = (-1)^{\operatorname{Card}(A)}\) . Let \(G\) be an abelian group written additively, and let \(\varphi\) and \(\psi\) be two maps from \(\mathfrak{F}\) to \(G\) . Show that the following properties are equivalent:
\end{enumerate}

\[
\varphi (\mathrm{A}) = \sum_ {\mathrm{B} \subset \mathrm{A}} \psi (\mathrm{B}) \quad \text {   for   all   } \mathrm{A} \in \mathfrak {F};\tag{i}
\]

\[
\psi (A) = \sum_ {B \subset A} \varepsilon (A - B) \varphi (B) \quad \text {   for   all   } A \in \mathfrak {F}.\tag{ii}
\]

\begin{enumerate}
\def\labelenumi{\arabic{enumi})}
\setcounter{enumi}{25}
\tightlist
\item
  Let \((W, S)\) be a Coxeter system, with S finite. For any subset H of W, denote by \(\mathrm{H}(t)\) the formal power series with integer coefficients defined by
\end{enumerate}

\[
\mathrm{H} (t) = \sum_ {w \in \mathrm{H}} t ^ {l (w)}.
\]

\begin{enumerate}
\def\labelenumi{\alph{enumi})}
\tightlist
\item
  Assume that \(\operatorname{Card}(S) = 2\) . Show that
\end{enumerate}

\[
\mathrm{W} (t) = \left(1 + t - t ^ {m} - t ^ {m + 1}\right) / (1 - t) \quad \text { if   } \mathrm{W} \text {   is   of   finite   order   } 2 m,
\]

\[
\mathrm{W} (t) = (1 + t) / (1 - t) \quad \text { if   } \mathrm{W} \text {   is   infinite. }
\]

\begin{enumerate}
\def\labelenumi{\alph{enumi})}
\setcounter{enumi}{1}
\tightlist
\item
  Assume that W is finite. Let \(w_0\) be the element of W of greatest length (Exerc. 22) and let \(m = l(w_0)\) . Show that
\end{enumerate}

\[
\mathrm{W} (t) = t ^ {m} \mathrm{W} (t ^ {- 1})
\]

(use Exerc. 22).

\begin{enumerate}
\def\labelenumi{\alph{enumi})}
\setcounter{enumi}{2}
\tightlist
\item
  Let X be a subset of S. Denote by \(A_{X}\) the set of \((X,\varnothing)\) -reduced elements of W (Exerc. 3) and \(W_{X}\) the subgroup of W generated by X. We know (Exerc. 3) that an element \(w \in W\) belongs to \(A_{X}\) if and only if \(l(xw) = l(w) + 1\) for all \(x \in X\) , that every element \(w \in W\) can be written uniquely in the form w = uv with \(u \in W_{X}\) and \(v \in A_{X}\) , and that in that case \(l(w) = l(u) + l(v)\) . Deduce the formula
\end{enumerate}

\[
\mathrm{W} (t) = \mathrm{W} _ {\mathrm{X}} (t). \mathrm{A} _ {\mathrm{X}} (t).
\]

\(d\) ) Retain the above notation and denote by \(B_{X}\) the set of \(w\in A_X\) such that \(l(sw) = l(w) - 1\) for all \(s\in S - X\) . Show that

\[
A _ {X} (t) = \sum_ {X \subset Y \subset S} B _ {Y} (t).
\]

Deduce that

\[
\mathrm{B} _ {\mathrm{X}} (t) = \sum_ {\mathrm{X} \subset \mathrm{Y} \subset \mathrm{S}} \varepsilon (\mathrm{Y} - \mathrm{X}) \mathrm{A} _ {\mathrm{Y}} (t) \quad \text { with } \quad \varepsilon (\mathrm{Z}) = (- 1) ^ {\operatorname{Card} (\mathrm{Z})}
\]

(use Exerc. 25).

\begin{enumerate}
\def\labelenumi{\alph{enumi})}
\setcounter{enumi}{4}
\tightlist
\item
  Assume that W is finite and define m and \(w_{0}\) as in b). Show that \(B_{\varnothing} = \{w_{0}\}\) and that
\end{enumerate}

\[
t ^ {m} = \sum_ {\mathrm{Y} \subset \mathrm{S}} \varepsilon (\mathrm{Y}) \frac {\mathrm{W} (t)}{\mathrm{W} _ {\mathrm{Y}} (t)}
\]

(use c) and d)).

\begin{enumerate}
\def\labelenumi{\alph{enumi})}
\setcounter{enumi}{5}
\tightlist
\item
  Assume that W is infinite. Show that \(B_{\varnothing} = \varnothing\) and that
\end{enumerate}

\[
0 = \sum_ {Y \subset S} \frac {\varepsilon (Y)}{W _ {Y} (t)}.
\]

\(g\) ) Show that the formal power series \(W(t)\) is a rational function of \(t\) (use \(f\) ) and argue by induction on \(\operatorname{Card}(\mathbf{S})\) ). Show that this rational function vanishes whenever \(t\) is a root of unity, and that \(1 / W(\infty)\) is an integer. Show that \(\frac{1}{W(t^{-1})} \in \mathbf{Z}[[t]]\) .

\exercisesection{Exercises for § 2}

\begin{enumerate}
\def\labelenumi{\arabic{enumi})}
\item
  Let G be a group, B and N two subgroups of G and S a subset of \(W = N/(B \cap N)\) . For all \(w \in W\) , put \(C(w) = BwB\) . Assume that conditions (T1) and (T2) of Def. 1 of no. 1 are satisfied, that whenever \(s \in S\) and \(w \in W\) at least one of the two relations \(C(s).C(w) = C(sw)\) and \(C(s).C(sw) = C(w)\) is satisfied, and that \(B \cup C(s)\) is a subgroup of G for all \(s \in S\) . Show that condition (T3) is satisfied. Moreover, if the index of B in \(B \cup C(s)\) is \(\geqslant 3\) , then (G, B, N, S) is a Tits system.
\item
  Let G be a group, B and N two subgroups of G and S a subset of \(W = N/(B \cap N)\) . Let Z be a normal subgroup of G contained in B. Let \(B'\) and \(N'\) be the canonical images of B and N in \(G' = G/Z\) . Show that the canonical map from N to \(N'\) defines an isomorphism from W to \(W' = N'/(B' \cap N')\) . Let \(S'\) be the image of S under this isomorphism. Show that \((G', B', N', S')\) is a Tits system if and only if \((G, B, N, S)\) is one.
\item
  Let G be a group, B a subgroup of G and \((\mathrm{C}(w))_{w\in\mathbb{W}}\) the family of double cosets of G with respect to B. Then B is called a Tits subgroup of G if there exists a subset S of W such that the following conditions are satisfied:
\end{enumerate}

\begin{enumerate}
\def\labelenumi{(\arabic{enumi})}
\item
  the union of the \(C(s)\) for \(s \in S\) generates \(G\) ;
\item
  for all \(s \in S\) , the set \(B \cup C(s)\) is a subgroup of \(G\) and the index of \(B\) in \(B \cup C(s)\) is \(\geqslant 3\) ;
\item
  for all \(s \in S\) and all \(w \in W\) , there exists an element \(w' \in W\) such that \(C(s).C(w) \subset C(w) \cup C(w')\) .
\end{enumerate}

From now on we assume that B is a Tits subgroup of G and that we are given a subset S of W satisfying conditions (1), (2) and (3).

\begin{enumerate}
\def\labelenumi{\alph{enumi})}
\item
  Show that \(C(s)^{-1} = C(s)\) and that \(C(s).C(s) = B \cup C(s)\) for all \(s \in S\) . Show that, for all \(s \in S\) and all \(w \in W\) , there exists an element \(w'' \in W\) such that \(C(w).C(s) \subset C(w) \cup C(w'')\) .
\item
  If \(w \in \mathbf{W}\) , the length of \(w\) , denoted by \(l(w)\) , is the smallest integer \(n \geqslant 0\) for which there exist \(s_1, \ldots, s_n \in \mathbf{S}\) with \(\mathrm{C}(w) \subset \mathrm{C}(s_1) \ldots \mathrm{C}(s_n)\) . Show that \(l(w)\) is finite for all \(w \in \mathbf{W}\) .
\end{enumerate}

Let \(u, v \in W\) with \(l(u) < l(v)\) and let \(s \in S\) . Show that, if \(C(v) \subset C(u).C(s)\) (resp. \(C(v) \subset C(s).C(u)\) ), then \(C(v) = C(u).C(s)\) (resp., \(C(v) = C(s).C(u)\) ) (argue by induction on the length of \(u\) . If \(C(v) \neq C(u).C(s)\) , then \(C(u).C(s) = C(u) \cup C(v)\) . By using the induction hypothesis, show that there exist \(t \in S\) and \(w \in W\) such that

\[
\mathrm{C} (u) = \mathrm{C} (t). \mathrm{C} (w) \quad \text { with } \quad l (w) = l (u) - 1.
\]

From the relations \(\mathrm{C}(v) \subset \mathrm{C}(u).\mathrm{C}(s) = \mathrm{C}(t).\mathrm{C}(w).\mathrm{C}(s)\) and \(\mathrm{C}(t).\mathrm{C}(w) = \mathrm{C}(u) \neq \mathrm{C}(v)\) , deduce the existence of an element \(w' \neq w\) such that

\(C(w') \subset C(w).C(s)\) and \(C(v) \subset C(t).C(w')\) , so \(l(w') \geqslant l(v) - 1 > l(u) - 1 = l(w)\) . The induction hypothesis now implies that

\[
\mathrm{C} (w ^ {\prime}) = \mathrm{C} (w). \mathrm{C} (s).
\]

Moreover,

\[
\begin{array}{c} \mathrm{C} (t). \mathrm{C} (u). \mathrm{C} (s) = \mathrm{C} (t). \mathrm{C} (t). \mathrm{C} (w). \mathrm{C} (s) = \mathrm{C} (t). \mathrm{C} (w). \mathrm{C} (s) \cup \mathrm{C} (w). \mathrm{C} (s) \\ = \mathrm{C} (u). \mathrm{C} (s) \cup \mathrm{C} (w ^ {\prime}) = \mathrm{C} (u) \cup \mathrm{C} (v) \cup \mathrm{C} (w ^ {\prime}) \end{array}
\]

and also

\[
\mathrm{C} (t). \mathrm{C} (u). \mathrm{C} (s) = \mathrm{C} (t). (\mathrm{C} (u) \cup \mathrm{C} (v)) \supset \mathrm{C} (t). \mathrm{C} (t). \mathrm{C} (w) \supset \mathrm{C} (w),
\]

which is a contradiction since \(w \neq u, v, w'\) .

\begin{enumerate}
\def\labelenumi{\alph{enumi})}
\setcounter{enumi}{2}
\tightlist
\item
  Show that, for all \(w \in \mathbf{W}\) and all \(s \in S\) , there exists a unique element, denoted by \(s.w\) (resp. \(w * s\) ), distinct from \(w\) and such that
\end{enumerate}

\[
\begin{array}{c} \mathrm{C} (s. w) \subset \mathrm{C} (s). \mathrm{C} (w) \subset \mathrm{C} (w) \cup \mathrm{C} (s. w) \\ (\text {resp.} \mathrm{C} (w * s) \subset \mathrm{C} (w). \mathrm{C} (s) \subset \mathrm{C} (w) \cup \mathrm{C} (w * s)). \end{array}
\]

(Show by induction on \(l(w)\) that \(C(s).C(w) \neq C(w)\) . For this, write

\[
\mathrm{C} (w) = \mathrm{C} (u). \mathrm{C} (t)
\]

with \(t \in S\) , \(u \in W\) and \(l(w) = l(u) + 1\) . If \(C(s).C(w) = C(w)\) , then

\[
\mathrm{C} (s). \mathrm{C} (u). \mathrm{C} (t) = \mathrm{C} (u). \mathrm{C} (t)
\]

and multiplying on the right by \(\mathrm{C}(t)\) , we obtain that \(\mathrm{C}(u) \cup \mathrm{C}(w) = \mathrm{C}(s).\mathrm{C}(u) \cup \mathrm{C}(w)\) . Since \(\mathrm{C}(u) \neq \mathrm{C}(s).\mathrm{C}(u)\) by the induction hypothesis, we have, by b),

\[
\mathrm{C} (s). \mathrm{C} (u) = \mathrm{C} (w),
\]

and hence

\[
\mathrm{C} (u) \subset \mathrm{C} (s). \mathrm{C} (s). \mathrm{C} (u) = \mathrm{C} (s). \mathrm{C} (w) = \mathrm{C} (w),
\]

which is absurd).

\begin{enumerate}
\def\labelenumi{\alph{enumi})}
\setcounter{enumi}{3}
\item
  Let \(s \in S\) . Show that the map \(p_s: w \mapsto s.w\) (resp. \(q_s: w \mapsto w * s\) ) is a permutation of \(W\) and that \(p_s^2 = \operatorname{Id}(\operatorname{resp.} q_s^2 = \operatorname{Id})\) . Show that \(p_s \circ q_t = q_t \circ p_s\) for all \(s, t \in S\) (study the product \(C(s).C(w).C(t)\) for \(w \in W\) and show that \((s.w) * t \in \{w, s.w, w * t, s.(w * t)\}\) ; show that \((s.w) * t \notin \{s.w, w * t\}\) and that if \((s.w) * t = w\) then \(s.w = w * t\) and \(w = s.(w * t)\) ).
\item
  Show that the group of permutations P (resp. Q) generated by the \(p_{s}\) (resp. \(q_{s}\) ) for \(s \in S\) acts simply-transitively on W (to show that P is transitive, use (2) and argue as in Lemma 1 of no. 1; to show that P is simply transitive, use d)). Deduce that W has a unique group structure such that the map \(p \mapsto p(e)\)
\end{enumerate}

(where \(e\) denotes the element of \(W\) such that \(B = C(e)\) ) is an isomorphism from \(P\) to \(W\) . The map \(q \mapsto q(e)\) is then also an isomorphism; moreover, \(s.w = w * s = sw\) for all \(s \in S\) and all \(w \in W\) and \(C(w)^{-1} = C(w^{-1})\) .

\begin{enumerate}
\def\labelenumi{\alph{enumi})}
\setcounter{enumi}{5}
\item
  Show that the pair \((W,S)\) is a Coxeter system and generalize the results of no. 4.
\item
  Let X be a subset of S and \(W_{X}\) the subgroup of W generated by X. Show that the union \(G_{X}\) of the \(C(w)\) for \(w \in W_{X}\) is a subgroup of G and that Th. 3 of no. 5 is still true. Show that B is a Tits subgroup of \(G_{X}\) . Generalize Props. 2 and 3 of no. 5, and Def. 2, Prop. 4 and Th. 4 of no. 6. Show that S is the set of \(w \in W\) such that \(B \cup C(w)\) is a subgroup of G distinct from B. The Coxeter system (W, S) and the group W (called the Weyl group of (G, B)) thus depend only on (G, B).
\item
  Let N be a subgroup of G such that \(B \cap N\) is normal in N and such that the intersection with N of every double coset \(C(w)\) with respect to B is a double coset with respect to \(B \cap N\) . Show that the group \(N/(B \cap N)\) can be identified with W and that \((G, B, N, S)\) is a Tits system.
\end{enumerate}

\begin{enumerate}
\def\labelenumi{\arabic{enumi})}
\setcounter{enumi}{3}
\tightlist
\item
  Let \((G, B, N, S)\) and \((G', B', N', S')\) be two Tits systems with \(G = G'\) and \(B = B'\) , and with Weyl groups W and \(W'\) . Let j be the bijection from W to \(W'\) defined by the relation
\end{enumerate}

\[
\mathrm{B} w \mathrm{B} = \mathrm{B} ^ {\prime} j (w) \mathrm{B} ^ {\prime}.
\]

Show that \(j\) is a group isomorphism between \(W\) and \(W'\) and that \(j(S) = S'\) .

\begin{enumerate}
\def\labelenumi{\arabic{enumi})}
\setcounter{enumi}{4}
\tightlist
\item
  Let \(\Sigma = (G, B, N, S)\) be a Tits system. Put \(T = B \cap N\) and let \(\tilde{N}\) be the normaliser of \(N\) .
\end{enumerate}

\begin{enumerate}
\def\labelenumi{\alph{enumi})}
\tightlist
\item
  Let \(b \in B \cap \hat{N}\) . Show that \(bnb^{-1}n^{-1} \in B \cap N\) for all \(n \in N\) (put \(bn = n'b\) and use Th. 1) and that b belongs to the intersection \(\tilde{T}\) of the conjugates \(nBn^{-1}\) for \(n \in N\) . Show that \(\tilde{T} \cap N = T\) .
\end{enumerate}

If \(\tilde{\mathrm{T}} = \mathrm{T}\) , \(\Sigma\) is said to be saturated.

\begin{enumerate}
\def\labelenumi{\alph{enumi})}
\setcounter{enumi}{1}
\item
  Put \(\tilde{N} = N.\tilde{T}\) . Show that \(\tilde{N}\) is a subgroup of G containing \(\tilde{T}\) as a normal subgroup and that \(\tilde{N} \cap B = \tilde{T}\) . Show that the injection of N into \(\tilde{N}\) defines an isomorphism from the Weyl group W of \(\Sigma\) to \(\tilde{N}/\tilde{T}\) .
\item
  Show that \((\mathbf{G},\mathbf{B},\tilde{\mathbf{N}},j(\mathbf{S}))\) is a saturated Tits system, which is said to be associated to \(\Sigma\) .
\end{enumerate}

\begin{enumerate}
\def\labelenumi{\arabic{enumi})}
\setcounter{enumi}{5}
\item
  We use the notation of no. 2 and let \(\mathbf{N}_0\) be the subgroup of \(\mathbf{N}\) consisting of the matrices all of whose entries are equal to 0 or 1. Show that \(\mathbf{B} \cap \mathbf{N}_0 = \mathbf{T} \cap \mathbf{N}_0 = \{1\}\) and that the canonical map \(j\) from \(\mathbf{N}_0\) to \(\mathbf{W} = \mathbf{N} / \mathbf{T}\) is an isomorphism. Put \(S_0 = j^{-1}(S)\) . Show that \((\mathbf{G}, \mathbf{B}, \mathbf{N}_0, S_0)\) is a Tits system and that \((\mathbf{G}, \mathbf{B}, \mathbf{N}, S)\) is the associated saturated Tits system.
\item
  Let G be a group acting on a set E. Then G is said to act doubly-transitively on E if, whenever \(x, y, x', y' \in E\) are such that \(x \neq y\) and \(x' \neq y'\) , there exists an element \(g \in G\) such that \(g.x = x'\) and \(g.y = y'\) .
\end{enumerate}

\begin{enumerate}
\def\labelenumi{\alph{enumi})}
\item
  Let \((G, B, N, S)\) be a Tits system whose Weyl group W is of order 2. Show that G acts doubly-transitively on G/B.
\item
  Let G be a group acting doubly transitively on a set E. Assume that \(\operatorname{Card}(E) \geqslant 3\) . Let \(e \in E\) and let B be the stabilizer of e. Let \(x \in E\) , with \(x \neq e\) , and let \(n \in G\) be such that \(n(e) = x\) and \(n(x) = e\) . Let N be the subgroup of G generated by n.~Let s be the canonical image of n in N/T. Show that \((G, B, N, \{s\})\) is a Tits system whose Weyl group is of order 2.
\end{enumerate}

\begin{enumerate}
\def\labelenumi{\arabic{enumi})}
\setcounter{enumi}{7}
\tightlist
\item
  Let \((\mathrm{G},\mathrm{B},\mathrm{N},\mathrm{S})\) be a Tits system; put \(T=B\cap N\) and W=N/T. Let \(\tilde{G}\) be a group containing G as a normal subgroup. Assume that for all \(h\in\tilde{G}\) , there exists \(g\in G\) such that \(hBh^{-1}=gBg^{-1}\) and \(hNh^{-1}=gNg^{-1}\) . Let \(\hat{B}\) (resp. \(\hat{N}\) ) be the normaliser of B (resp. N) in \(\tilde{G}\) ; put \(\Gamma=\hat{B}\cap\hat{N}\) , \(\hat{N}=\Gamma.N\) and \(\tilde{T}=\tilde{N}\cap B\) .
\end{enumerate}

\begin{enumerate}
\def\labelenumi{\alph{enumi})}
\item
  Show that \(\tilde{\mathbf{G}} = \Gamma .\mathbf{G}\) , \(\hat{\mathbf{B}} = \Gamma .\mathbf{B}\) , \(\Gamma \cap \mathrm{B} = \Gamma \cap \mathrm{G}\) and \(\tilde{\mathbf{T}} = (\Gamma \cap \mathbf{B}).\mathbf{T}\) . The groups \(\Omega = \Gamma /(\Gamma \cap \mathrm{B})\) , \(\tilde{\mathbf{G}} /\mathbf{G}\) and \(\hat{\mathbf{B}} /\mathbf{B}\) are thus canonically isomorphic. If \(\varPhi \subset \varOmega\) and if \(\mathbf{H}\) is a subgroup of \(\mathbf{G}\) containing \(\Gamma \cap \mathrm{B}\) , we denote by \(\varPhi \mathrm{H}\) the union of the subsets \(\varphi \mathrm{H}\) for \(\varphi \in \varPhi\) .
\item
  Show that \(\tilde{T}\) is normal in \(\tilde{N}\) (to show that \(n\gamma n^{-1} \in \tilde{N}\) for \(n \in \tilde{N}\) and \(\gamma \in \Gamma \cap B\) , use Exerc. 5 a)), that \(N \cap \tilde{T} = T\) and that \(\Gamma \cap \tilde{T} = \Gamma \cap B\) . The injection of N (resp. \(\Gamma\) ) into \(\tilde{N}\) thus allows W (resp. \(\Omega\) ) to be identified with a subgroup of \(\tilde{W} = \tilde{N}/\tilde{T}\) . Show that \(\Omega\) normalises S and that \(\tilde{W}\) is the semi-direct product of \(\Omega\) and W.
\item
  Show that, for all \(s \in S\) and all \(u \in \tilde{W}\) ,
\end{enumerate}

\[
\mathrm{BsBuB} \subset (\mathrm{BuB}) \cup (\mathrm{BsuB}).
\]

\begin{enumerate}
\def\labelenumi{\alph{enumi})}
\setcounter{enumi}{3}
\item
  Show that the map \(u \mapsto BuB\) is a bijection from \(\tilde{W}\) to \(B\backslash\tilde{G}/B\) (use Th. 1 and the fact that \(\Gamma\) normalises B).
\item
  Let \(\mathfrak{B}\) be the set of pairs \((\varPhi, X)\) , where \(\varPhi\) is a subgroup of \(\Omega\) and \(X\) is a subset of \(X\) normalised by \(\varPhi\) . Put \(G_{(\varPhi, X)} = \varPhi G_X = B \varPhi W_X B\) (with the notation of no. 5). Show that the map \((\varPhi, X) \mapsto G_{(\varPhi, X)}\) is a bijection from \(\mathfrak{B}\) to the set of subgroups of \(\tilde{G}\) containing \(B\) . Generalize assertions \(b)\) and \(c)\) of Th. 3 and Prop. 2 of no. 5.
\end{enumerate}

Show that the normaliser of \(\mathrm{G}_{(\varPhi,\mathrm{X})}\) in \(\tilde{G}\) is the subgroup \(\mathrm{G}_{(\varPhi',\mathrm{X})}\) , where \(\varPhi'\) is the set of elements of \(\Omega\) normalizing both \(\varPhi\) and X.

\begin{enumerate}
\def\labelenumi{\alph{enumi})}
\setcounter{enumi}{5}
\tightlist
\item
  Show that \(G_{(\Phi, \mathsf{X})}\) is a maximal subgroup of \(\tilde{\mathbf{G}}\) if and only if one of the following two conditions is satisfied:
\end{enumerate}

\begin{enumerate}
\def\labelenumi{(\roman{enumi})}
\item
  \(X = S\) and \(\varPhi\) is a maximal subgroup of \(\varOmega\) ;
\item
  \(\Phi = \Omega\) and \(\Phi\) acts transitively on \(S - X\) (which is non-empty).
\end{enumerate}

Show that \(G_{(\Phi,X)}\) is a maximal subgroup in the set of subgroups of \(\tilde{G}\) not containing G if and only if

\begin{enumerate}
\def\labelenumi{(\roman{enumi})}
\setcounter{enumi}{2}
\tightlist
\item
  \(X \neq S\) , \(\Phi\) is the normaliser of \(X\) in \(\Omega\) and acts transitively on \(S - X\) .
\end{enumerate}

\(g)\) Let \(\Phi\) be a normal subgroup of \(\Omega\) and put \(G' = \Phi G\) , \(B' = \Phi B\) , \(N' = \Phi N\) and \(T' = B' \cap N'\) . Show that \(T' = \Phi \tilde{T}\) and that \(T'\) is normal in \(N'\) if and only if every element of \(\Phi\) commutes with every element of \(W\) . Show that \(G'\) is then normal in \(\tilde{G}\) , that the injection of \(N\) into \(N'\) defines an isomorphism \(j\) from \(W\) to \(W' = N'/T'\) and that \((G', B', N', S')\) (with \(S' = j(S)\) ) is a Tits system.

\begin{enumerate}
\def\labelenumi{\arabic{enumi})}
\setcounter{enumi}{8}
\tightlist
\item
  Let \((G, B, N, S)\) be a Tits system and X, Y, Z three subsets of S. Show that
\end{enumerate}

\[
\mathrm{G} _ {\mathrm{X}} \cap (\mathrm{G} _ {\mathrm{Y}}. \mathrm{G} _ {\mathrm{Z}}) = (\mathrm{G} _ {\mathrm{X}} \cap \mathrm{G} _ {\mathrm{Y}}). (\mathrm{G} _ {\mathrm{X}} \cap \mathrm{G} _ {\mathrm{Z}})
\]

(use Exerc. 1 of §1 and Prop. 2 of no. 3).

\begin{enumerate}
\def\labelenumi{\arabic{enumi})}
\setcounter{enumi}{9}
\tightlist
\item
  Let G be a group and B a Tits subgroup of G. We use the notation of Exerc. 3. For \(s \in S\) , denote by \(G^{(s)}\) the subgroup \(G_{S - \{s\}}\) (Exerc. 3 g)). Let I be the set of subsets of G of the form \(gG^{(s)}\) (for \(g \in G\) and \(s \in S\) ), and C the set of subsets of I of the form \(C_g = \{gG^{(s)} \mid s \in S\}\) for \(g \in G\) . The \(C_g\) are called the chambers of I (§ 1, Exerc. 15). Let G act on I by left translations.
\end{enumerate}

\begin{enumerate}
\def\labelenumi{\alph{enumi})}
\item
  Let \(\mathfrak{F}\) be the set of facets of I (that is, the subsets of the chambers, cf.~§1, Exerc. 15). Let \(F \in \mathfrak{F}\) . Show that there exist a unique subset X of S and an element \(g \in G\) such that \(gG_{X} = \bigcap_{a \in F} a\) ; F is said to be of type X. Show that F is then the set of \(gG^{(s)}\) for \(s \in S - X\) and is of codimension \(\operatorname{Card}(X)\) in any chamber containing it. Show that the map \(j: F \mapsto \bigcap_{a \in F} a\) is a strictly decreasing bijection, compatible with left translations, from \(\mathfrak{F}\) to the set of subsets of G of the form \(gG_{X}\) for \(g \in G\) and \(X \subset S\) . Show that, if \(X \subset Y \subset S\) , every facet of type X contains a unique facet of type Y.
\item
  Show that G acts transitively on the set C of chambers and that the stabilizer of the chamber \(C_{g}\) ( \(g \in G\) ) is equal to \(gBg^{-1}\) ; the map \(g \mapsto C_{g}\) thus defines a bijection from G/B to C.
\item
  Show that two chambers \(C_{g}\) and \(C_{g'}\) ( \(g, g' \in G\) ) are neighbouring (§1, Exerc. 15) if and only if there exists \(s \in S\) such that \(g' \in g(B \cup BsB)\) .
\end{enumerate}

\(d)\) Let \(C_1, \ldots, C_n\) be chambers of I and put \(C_0 = C_e\) . Show that the following conditions are equivalent:

\begin{enumerate}
\def\labelenumi{(\roman{enumi})}
\item
  the sequence \(\Gamma = (\mathrm{C}_0, \mathrm{C}_1, \ldots, \mathrm{C}_n)\) is an injective gallery;
\item
  there exist a sequence \(\mathbf{s} = (s_1, \ldots, s_n)\) of elements of S and a sequence \((b_1, \ldots, b_n)\) of elements of B such that \(C_j = b_1 \bar{s}_1 b_2 \bar{s}_2 \ldots b_j \bar{s}_j (C_0)\) (where \(\bar{s}_j\) denotes a given element of the double coset \(Bs_jB\) ) for \(1 \leqslant j \leqslant n\) .
\end{enumerate}

Show that, if these conditions are satisfied, the sequence s is unique: it is then called the type of \(\Gamma\) and denoted by \(\mathbf{s}(\Gamma)\) . Show that an injective gallery is minimal if and only if its type is a reduced decomposition. Show that the following conditions are satisfied:

(WI 1) For any two chambers C and \(\mathbf{C}'\) of I, there exists a unique element \(t(\mathbf{C},\mathbf{C}')\) of W such that the set of types of minimal galleries with extremities C and \(\mathbf{C}'\) is the set of reduced decompositions of \(t(\mathbf{C},\mathbf{C}')\) .

(WI 2) For any chamber C, the map \(\mathrm{C}' \mapsto t(\mathrm{C}, \mathrm{C}')\) from the set of chambers of I to W is surjective.

\begin{enumerate}
\def\labelenumi{\alph{enumi})}
\setcounter{enumi}{4}
\item
  Show that two minimal galleries of the same type and with the same extremities are identical. (Reduce to proving that, if \((s_{1},\ldots,s_{n})\) is a reduced decomposition and if \(b_{1},\ldots,b_{n},b_{1}^{\prime},\ldots,b_{n}^{\prime}\) are elements of B such that \(b_{1}\bar{s}_{1}\ldots b_{n}\bar{s}_{n}\in b_{1}^{\prime}\bar{s}_{1}\ldots b_{n}^{\prime}\bar{s}_{n}B\) , then \(b_{1}\bar{s}_{1}\in b_{1}^{\prime}\bar{s}_{1}B\) . For this, remark that if \(\bar{s}_{1}^{-1}b_{1}^{-1}b_{1}^{\prime}\bar{s}_{1}\notin B\) , then this element would belong to \(Bs_{1}B\) and we would have \(b_{2}\bar{s}_{2}\ldots b_{n}\bar{s}_{n}\in b\bar{s}_{1}b^{\prime}b_{2}^{\prime}\bar{s}_{2}\ldots b_{n}^{\prime}\bar{s}_{n}B\) with \(b,b^{\prime}\in B\) , contradicting Cor. 1 of Th. 2 of no. 4.)
\item
  Show that I equipped with C is a building, said to be associated to the pair (G, B). Show that there exists a unique numbering (§ 1, Exerc. 20) of I for which the type of a facet is that defined in a).
\item
  Show that I is spacious, namely that every panel of I is contained in at least three chambers (cf.~§ 1, Exerc. 24). Show that the following condition is satisfied:
\end{enumerate}

\begin{enumerate}
\def\labelenumi{(\Alph{enumi})}
\setcounter{enumi}{6}
\tightlist
\item
  Given a panel F, a chamber C, a gallery \(\Gamma = (\mathrm{C}_0, \ldots, \mathrm{C}_n)\) such that \(\mathrm{F} \subset \mathrm{C}_n\) and of smallest possible length, and chambers \(\mathrm{C}'\) and \(\mathrm{C}''\) containing F and distinct from \(\mathrm{C}_n\) , there exists an element \(g \in \mathrm{G}\) such that \(g(\mathrm{C}_i) = \mathrm{C}_i\) for \(0 \leqslant i \leqslant n\) and \(g(\mathrm{C}') = \mathrm{C}''\) .
\end{enumerate}

(Reduce to the case \(C_{0} = C\) ; let \(u \in S\) be such that F is of type \(S - \{u\}\) and let s be the type of \(\Gamma\) . Show that \((\mathbf{s}, u)\) is a reduced decomposition. Take \(h \in G\) such that \(C_{n} = h(\mathrm{C})\) ; there exist \(b'\) and \(b'' \in B\) such that \(C' = hb'u(\mathrm{C})\) and \(C'' = hb''u(\mathrm{C})\) . Then use Cor. 1 of Th. 2 of no. 4 generalized to the case of a Tits subgroup (cf.~Exerc. 3 f)) to show that there exists \(b \in B\) such that \(bhb'uB = hb''uB\) ; then \(b \in hBuBu^{-1}Bh^{-1} \subset (hBh^{-1}) \cup (hBuBh^{-1})\) . If \(b \in hBuBh^{-1}\) , we would have BhB \(\subset BhBuB = BhuB\) (loc. cit.), which is impossible. Hence \(b(C_{n}) = C_{n}\) and e) implies that \(b(C_{i}) = C_{i}\) for all i.)

\begin{enumerate}
\def\labelenumi{\arabic{enumi})}
\setcounter{enumi}{10}
\tightlist
\item
  Let \((\mathbf{W},\mathbf{S})\) be a Coxeter system and I a building numbered by S (§ 1, Exerc. 20). If \(\Gamma = (\mathbf{C}_0,\dots ,\mathbf{C}_n)\) is an injective gallery, the type of \(\Gamma\) , denoted by \(s(\Gamma)\) , is the sequence \((s_1,\ldots ,s_n)\) of elements of S such that the panel \(\mathbf{C}_{i - 1}\cap \mathbf{C}_i\) is of type \(\mathbf{S} - \{s_i\}\) (for \(1\leqslant i\leqslant n\) ). If conditions (WI 1) and (WI 2) of Exerc. 10 are satisfied, I is called a \((\mathbf{W},\mathbf{S})\) -building.
\end{enumerate}

Let I be a spacious (W, S)-building (cf.~Exerc. 10 g)) and let G be a group of allowed automorphisms of I satisfying the following condition:

(G0) Given three distinct chambers C, C' and C'' containing the same panel, there exists an element \(g \in G\) such that \(g(\mathrm{C}) = \mathrm{C}\) and \(g(\mathrm{C}') = \mathrm{C}'\) .

Choose a chamber C of I and denote by B the stabilizer of C in G.

\begin{enumerate}
\def\labelenumi{\alph{enumi})}
\item
  Show that G acts transitively on the set of chambers of I.
\item
  Let \(C'\) and \(C''\) be two chambers. Show that \(t(\mathrm{C},\mathrm{C}') = t(\mathrm{C},\mathrm{C}'')\) if and only if there exists \(b \in B\) such that \(\mathrm{C}'' = b(\mathrm{C}')\) (if \(t(\mathrm{C},\mathrm{C}') = t(\mathrm{C},\mathrm{C}'') = w\) , consider a reduced decomposition s of w and a minimal gallery \(\Gamma'\) (resp. \(\Gamma''\) ) with extremities C and \(C'\) (resp. \(C''\) ) and of type s. Argue by induction on \(l(w)\) using \((\mathrm{G}_{0})\) and a)). Deduce that there exists a bijection \(w \mapsto \mathrm{B}(w)\) from W to B\textbackslash G/B.
\item
  Let F be the facet of C of type S- \(\{s\}\) . Show that the stabilizer of F in G is equal to B \(\cup\) B(s) (if g(F) = F, then t(C, g(C)) = 1 or s). Show that B \(\subset\) B(s)B(s) (use (G0)).
\item
  Show that B is a Tits subgroup of G and that the Coxeter system of (G,B) is canonically isomorphic to (W,S). (Let \(w \in W\) and \(s \in S\) be such that \(l_{\mathbb{S}}(sw) = l_{\mathbb{S}}(w) + 1\) . Let \(g \in \mathrm{B}(w)\) and \(u \in \mathrm{B}(s)\) ; let \(\mathbf{s} = (s_1, \ldots, s_n)\) be a reduced decomposition of \(w\) , ( \(C = C_0, C_1, \ldots, C_n = g(C)\) ) a minimal gallery of type s and choose an element \(\bar{s}_i \in \mathrm{B}(s_i)\) ; show, by using \((G_0)\) , that there exist elements \(b_i \in B\) such that \(C_i = b_1 \bar{s}_1 \ldots b_i \bar{s}_i(C)\) . Put \(C_i' = u(C_{i-1})\) and show that the gallery \((C, u(C), u(C_1), \ldots, u(C_n))\) is of type \((s, s)\) , and is therefore minimal. Deduce that \(ug \in \mathrm{B}(sw)\) and that \(B(s) B(w) = B(sw)\) . If now \(l(sw) = l(w) - 1\) , put \(w' = sw\) : then \(B(s) B(w') = B(w)\) and hence
\end{enumerate}

\[
\mathrm{B} (s) \mathrm{B} (w) = \mathrm{B} (s) \mathrm{B} (s) \mathrm{B} \left(w ^ {\prime}\right) \subset \mathrm{B} \left(w ^ {\prime}\right) \cup \left(\mathrm{B} (s) \mathrm{B} \left(w ^ {\prime}\right)\right) = \mathrm{B} (s w) \cup \mathrm{B} (w);
\]

finally, since \(\mathrm{B} \subset \mathrm{B}(s)\mathrm{B}(s)\) , we also have \(\mathrm{B}(w') = \mathrm{B}(sw) \subset \mathrm{B}(s)\mathrm{B}(w)\) .

\begin{enumerate}
\def\labelenumi{\alph{enumi})}
\setcounter{enumi}{4}
\tightlist
\item
  Show that there exists a unique isomorphism from the building associated to (G, B) (Exerc. 10) to I, compatible with the action of G and taking the canonical chamber \(C_{e}\) to C.
\end{enumerate}

\begin{enumerate}
\def\labelenumi{\arabic{enumi})}
\setcounter{enumi}{11}
\tightlist
\item
  Let \((\mathbf{G},\mathbf{B},\mathbf{N},\mathbf{S})\) be a Tits system, \(\mathrm{W} = \mathrm{N} / (\mathrm{B}\cap \mathrm{N})\) its Weyl group, I the \((\mathrm{W},\mathrm{S})\) -building associated to \((\mathbf{G},\mathbf{B})\) (Exerc. 10) and \(\mathbf{C} = \mathbf{C}_e\) the canonical chamber of I. Let \(\mathbf{A}_0\) be the apartment associated to the Coxeter system \((\mathrm{W},\mathrm{S})\) (§ 1, Exerc. 16). For all \(g\in \mathbf{G}\) , let \(\varphi_g\) be the map from \(\mathbf{A}_0\) to I that takes a point \(w\mathrm{W}^{(s)}\) of \(\mathbf{A}_0\) ( \(w\in \mathrm{W}, s\in \mathrm{S}\) ) to the point \(gw\mathrm{G}^{(s)}\) of I.
\end{enumerate}

Show that, for all \(g \in G\) , the map \(\varphi_{g}\) is an isomorphism of numbered buildings from \(A_{0}\) to a subset of I that is the union of chambers \(gn(C_{e})\) for \(n \in N\) . Show that I, equipped with the set U of \(\varphi_{g}(A_{0})\) for \(g \in G\) , is a structured building (§1, Exerc. 24) (to prove (SB 2), remark that, if \(g', g'' \in G\) , there exist \(b', b'' \in B\) and \(n \in N\) such that \(g'^{-1}g'' = b'nb''\) ; then put \(g = g'b'n\) and show that \(g'(C)\) and \(g''(C)\) are contained in \(\varphi_{g}(A_{0})\) . To prove (SB 3), reduce to the case of two apartments \(A' = \varphi_{e}(A_{0})\) and

\(A'' = \varphi_b(A_0)\) , with \(b \in B\) , and show that the map \(a \mapsto b(a)\) leaves the points of \(A' \cap A''\) fixed by using Prop. 2 of no. 5). Show that the Coxeter system (W, S) and the numbering of I are adapted to (I, U) (§1, Exerc. 24 e)) and that the set \(\mathfrak{I}\) of allowed isomorphisms from \(A_0\) to the various elements of \(\mathfrak{U}\) is the set of \(\varphi_g\) for \(g \in G\) .

\begin{enumerate}
\def\labelenumi{\arabic{enumi})}
\setcounter{enumi}{12}
\tightlist
\item
  We use the notation of Exerc. 24 of §1: (I, U) is a spacious structured building equipped with a Coxeter system (W, S) and an adapted numbering and I is the set of allowed isomorphisms from the apartment \(A_{0}\) associated to (W, S) to the various apartments of (I, U). Further, let G be a group of allowed automorphisms of I, preserving E: the group G then acts on I and we assume that G acts transitively on I.
\end{enumerate}

Denote by C a chamber of I, A an apartment of U containing C, \(\varphi\) an allowed isomorphism from \(A_{0}\) to A taking the canonical chamber \(C_{e}\) of \(A_{0}\) to C, B the stabilizer of C in G and N the stabilizer of A in G.

\begin{enumerate}
\def\labelenumi{\alph{enumi})}
\item
  Show that, if \(A'\) and \(A''\) are two apartments belonging to U, and containing the same chamber, there exists \(g \in G\) such that \(g(A') = A''\) and that \(g(a) = a\) for all \(a \in A' \cap A''\) . Show that G acts transitively on the set of pairs (A, C) where \(A \in U\) and C is a chamber of A.
\item
  Show that the map \(n \mapsto \varphi^{-1} \circ n \circ \varphi\) is a surjective homomorphism from N to W, with kernel \(B \cap N\) . We can thus identify \(N/(B \cap N)\) with W.
\item
  Show that conditions (WI 1) and (WI 2) of Exerc. 10 d) are satisfied (use the following fact: if an apartment of I contains two chambers C' and C'', it contains every minimal gallery with extremities C' and C'' (§ 1, Exerc. 24 b))).
\item
  Show that condition (G) of Exerc. 10 g) (and a fortiori condition (G₀) of Exerc. 11) is satisfied (with the notation of (G), consider an apartment A' (resp. A'') of I containing C₀ and C' (resp. C''); show by using Exerc. 24 b) of §1 that Γ ⊂ A'' ∩ A'' and use a)).
\item
  Show that \((G, B, N, S)\) is a Tits system and that \((I, \mathfrak{U})\) is canonically isomorphic to the associated numbered structured building (Exerc. 12).
\end{enumerate}

\begin{enumerate}
\def\labelenumi{\arabic{enumi})}
\setcounter{enumi}{13}
\item
  Let \((\mathbf{G},\mathbf{B},\mathbf{N},\mathbf{S})\) be a Tits system and \((\mathbf{I},\mathfrak{U})\) the associated numbered structured building (Exerc. 12). Show that G, considered as a group of automorphisms of I, satisfies the conditions of Exerc. 13. With the notation of Exerc. 12, put \(\mathrm{A}=\varphi_{e}(\mathrm{A}_{0})\) and \(\mathrm{C}=\varphi_{e}(\mathrm{C}_{e})\) ; show that B is the stabilizer of C in G. Let \(\tilde{N}\) be the stabilizer of A in G; show that \(\tilde{\mathrm{N}}/(\mathrm{B}\cap\tilde{\mathrm{N}})\) can be identified with W and that \((\mathrm{G},\mathrm{B},\tilde{\mathrm{N}},\mathrm{S})\) is the saturated Tits system associated to \((\mathrm{G},\mathrm{B},\mathrm{N},\mathrm{S})\) (Exerc. 5).
\item
  We use the hypotheses and notation of Exerc. 10. Assume moreover that the Weyl group \(\mathbf{W}\) of \((\mathbf{G},\mathbf{B})\) is finite and denote by \(w_0\) the longest element of \(\mathbf{W}\) (§1, Exerc. 22). Two chambers \(\mathbf{C}\) and \(\mathbf{C}'\) are said to be opposed if \(t(\mathbf{C},\mathbf{C}') = w_0\) .
\end{enumerate}

\begin{enumerate}
\def\labelenumi{\alph{enumi})}
\item
  Show that if C and C' are opposed, so are C' and C. Show that there exists a chamber opposed to any given chamber C. Show that the stabilizer of a chamber C in G acts transitively on the set of chambers opposed to C.
\item
  Let C and \(C'\) be two opposed chambers. Show that for all \(w \in W\) , there exists a unique chamber \(C_w\) having the following property: if \((s_1, \ldots, s_k)\) (resp. \((s'_1, \ldots, s'_h)\) ) is a reduced decomposition of w (resp. of \(w' = w_0 w^{-1}\) ), there exists a minimal gallery \((C_0 = C, C_1, \ldots, C_n = C')\) of type \((s_1, \ldots, s_k, s'_1, \ldots, s'_h)\) (with \(n = h + k\) ) such that \(C_w = C_k\) (use Exerc. 22 of §1 and Exerc. 10, d) and e)). Show that \(C_w\) and \(C_{ww_0}\) are opposed.
\item
  Let \(\mathfrak{M}\) be the set of pairs of opposed chambers. For \(m = (C, C') \in \mathfrak{M}\) , let \(A_m\) be the union of the chambers \(C_w\) constructed above. Show that I equipped with the set \(\mathfrak{U}\) of \(A_m\) for \(m \in \mathfrak{M}\) is a structured building ( \(\S 1\) , Exerc. 24) and that the Coxeter system (W, S) and the numbering of I are adapted to (I, \(\mathfrak{U}\) ); define a canonical bijection from \(\mathfrak{M}\) to the set denoted by \(\mathfrak{I}\) in Exerc. 24 of \(\S 1\) . We identify these two sets.
\item
  Let \(m = (C, C') \in \mathfrak{M}\) with \(C = C_e\) , and let N be the stabilizer of \(A_m\) in G. Show that \(N/(B \cap N)\) can be identified with W and that \((G, B, N, S)\) is a saturated Tits system (use Exerc. 13).
\end{enumerate}

\begin{enumerate}
\def\labelenumi{\arabic{enumi})}
\setcounter{enumi}{15}
\tightlist
\item
  We retain the hypotheses and notation of Exerc. 15.
\end{enumerate}

\begin{enumerate}
\def\labelenumi{\alph{enumi})}
\item
  Let C and \(C'\) be two chambers. Show that there exists a chamber \(C''\) opposed both to C and to \(C'\) (take \(C''\) opposed to C such that \(t(\mathrm{C},\mathrm{C}')\) is of greatest possible length; if \(t(\mathrm{C}',\mathrm{C}'')\neq w_{0}\) , there exists a neighbouring chamber \(C_{1}\) of \(C''\) such that \(l(t(\mathrm{C}',\mathrm{C}_{1}))>l(t(\mathrm{C}',\mathrm{C}''))\) ; show by using condition (G) of Exerc. 10 that we can assume that \(C_{1}\not\subset A_{(C,C'')}\) and that \(C_{1}\) is then opposed to C).
\item
  Let \(a \in U\) ; show that there exists a unique involutive automorphism \(j_{A}\) (not necessarily allowed) that takes every chamber of A to the opposed chamber (use Exerc. 22 c) of §1). Let F and \(F'\) be two facets of I; show that if \(F' = j_{A}(F)\) for some \(A \in U\) containing F and \(F'\) , then the same is true for all \(A \in U\) containing F and \(F'\) (if \(F, F' \in A \cap A'\) , with \(A, A' \in U\) , consider a chamber C (resp. \(C'\) ) of A (resp. \(A'\) ) containing F (resp. \(F'\) ) and an \(A'' \in U\) containing C and \(C'\) and use (SB 3)). Then F and \(F'\) are said to be opposed. Show that two facets have a common opposed facet if and only if they are of the same type T and that a facet opposed to a facet of type T is of type \(w_{0}Tw_{0}^{-1}\) .
\item
  Let \(A_{0}\) be the apartment associated to a Coxeter system (W, S) (§1, Exerc. 16) and let J be the set of allowed isomorphisms from \(A_{0}\) to the various elements of U. If \(\alpha\) is a half of \(A_{0}\) , with wall L (§1, Exerc. 16) and if \(\varphi \in J\) , we say that \(\varphi(\alpha)\) is a semi-apartment of I, with wall \(\varphi(L)\) . Show that \(\varphi(L)\) then depends only on \(\varphi(\alpha)\) and not on the pair \((\varphi, \alpha)\) . Let \(D_{1}\) and \(D_{2}\) be two distinct semi-apartments, with the same wall L: show that there exist \(\varphi \in J\) and a wall \(L_{0}\) of \(A_{0}\) such that \(L = \varphi(L_{0})\) and that the \(D_{i}\) are the images under \(\varphi\) of the two halves of \(A_{0}\) determined by \(L_{0}\) (choose a panel F contained in L and two chambers \(C_{1}\) and \(C_{2}\) , one contained in \(D_{1}\) and the other in \(D_{2}\) and such that \(C_{1}\) contains F and \(C_{2}\) contains the panel opposed to F in \(D_{1}\) (which is also opposed to F in \(D_{2}\) ) and consider the apartment of I containing \(C_{1}\) and \(C_{2}\) ).
\end{enumerate}

\begin{enumerate}
\def\labelenumi{\arabic{enumi})}
\setcounter{enumi}{16}
\tightlist
\item
  We retain the hypotheses and notation of Exerc. 15 and 16. Choose an element \(\varphi \in \mathfrak{I}\) taking the canonical chamber \(C\) of \(A_0\) ( \(\S 1\) , Exerc. 16) to the canonical chamber \(C_e\) of \(I\) and, as in Exerc. 15 \(d\) ), denote by \(N\) the stabilizer of the apartment \(\varphi(A_0) \in \mathfrak{U}\) in the group \(G\) . For any subset \(D\) of \(A_0\) , denote by \(B_D\) the subgroup of \(G\) leaving fixed all the points of \(\varphi(D)\) : we have \(B = B_C\) and \(B \cap N = B_{A_0}\) .
\end{enumerate}

\begin{enumerate}
\def\labelenumi{\alph{enumi})}
\item
  Let \(\alpha\) be a half of \(A_0\) containing C. Show that \(B_{\alpha}\) acts transitively on the semi-apartments of I distinct from \(\varphi(\alpha)\) whose wall is the wall L of \(\varphi(\alpha)\) (let X be such a semi-apartment, F a panel contained in L and \(C'\) a chamber of \(\varphi(\alpha)\) containing F; show, by using Exerc. 16 c) and Exerc. 13 a), that there exists \(g \in G\) such that \(g(X \cup \varphi(\alpha)) = \varphi(A_0)\) and \(g(C') = C'\) ; show that \(g(F) = F\) and deduce from Exerc. 22 c) of §1 that \(g \in B_{\alpha}\) . Deduce that \(B_L\) acts doubly transitively on the semi-apartments with wall L and that \((B_L, B_{\alpha}, B_L \cap N)\) is a Tits system whose Weyl group is of order 2 (cf.~Exerc. 7).
\item
  Let \(D_{1}\) and \(D_{2}\) be two convex subsets of \(A_{0}\) such that \(C \subset D_{1} \subset D_{2}\) and that there exists a unique half \(\alpha\) of \(A_{0}\) with \(D_{1} \subset \alpha\) and \(D_{2} \not\subset \alpha\) : then \(D_{1} = D_{2} \cap \alpha\) (§1, Exerc. 21 c)). Show that \(B_{D_{1}} = B_{\alpha}B_{D_{2}}\) and that \(B_{\alpha} \cap B_{D_{2}} = B \cap N\) (by considering a gallery of smallest possible length having one extremity equal to C and the other containing a point \(a \in D_{2} - D_{1}\) , show that there exist two neighbouring chambers \(C_{1}'\) and \(C_{2}'\) of \(A_{0}\) such that \(C_{1}' \subset D_{1}\) , \(C_{2}' \subset D_{2}\) , \(C_{2}' \not\subset D_{1}\) and \(C_{1}' \cap C_{2}'\) contained in the wall \(L'\) of \(\alpha\) . Put \(C_{i} = \varphi(C_{i}')\) , \(F = \varphi(C_{1}' \cap C_{2}')\) and \(L = \varphi(L')\) . Show that the convex hull of \(D_{1} \cup C_{2}'\) is equal to \(D_{2}\) . Let \(b \in B_{D_{1}}\) : then \(b(C_{1}) = C_{1}\) , so \(b(F) = F\) and \(b(C_{2}) \neq C_{1}\) . Deduce that the convex hull of \(b(C_{2}) \cup L\) is a semi-apartment X distinct from \(\varphi(\alpha)\) and with wall L, and that there exists \(b' \in B_{\alpha}\) such that
\end{enumerate}

\[
b ^ {\prime} (\varphi (A _ {0})) = X \cup \varphi (\alpha).
\]

Show that \(b'b(a) = a\) for all \(a \in \varphi(\mathrm{D}_1 \cup \mathrm{C}_2')\) and that \(b'b \in \mathrm{B}_{\mathrm{D}_2}\) .

\begin{enumerate}
\def\labelenumi{\alph{enumi})}
\setcounter{enumi}{2}
\tightlist
\item
  Let \((\alpha_{i})_{1\leqslant i\leqslant q}\) be the halves of \(A_0\) containing C, numbered so that the map
\end{enumerate}

\[
j \mapsto \bigcap_ {1 \leqslant i \leqslant j} \alpha_ {i}
\]

is strictly decreasing (cf.~§1, Exerc. 17). Show that \(\mathbf{B} = \mathbf{B}_{\alpha_1} \ldots \mathbf{B}_{\alpha_q}\) and that if \(b_i, b_i' \in \mathbf{B}_{\alpha_i}\) with \(b_1 \ldots b_q = b_1' \ldots b_q'\) , then \(b_i' \in b_i(\mathbf{B} \cap \mathbf{N})\) for \(1 \leqslant i \leqslant q\) .

\begin{enumerate}
\def\labelenumi{(\arabic{enumi})}
\setcounter{enumi}{17}
\tightlist
\item
  We use the hypotheses and notation of no. 2. Let \(V_i\) be the vector subspace of \(k^n\) generated by \(e_1, \ldots, e_i\) (for \(1 \leqslant i \leqslant n - 1\) ).
\end{enumerate}

\begin{enumerate}
\def\labelenumi{\alph{enumi})}
\item
  Show that, for any subset X of S, the subgroup \(G_{X}\) (no. 5) consists of the elements \(g \in G\) such that \(g(V_{i}) = V_{i}\) for all i such that \(s_{i} \notin X\) .
\item
  Let I be the building associated to a Tits system (G, B, N, S) (Exerc. 10 and 12). Show that the map j that associates to the point \(g\mathbf{G}^{(i)}\) of I (where \(\mathbf{G}^{(i)}\) denotes the subgroup \(G_{S-\{s_i\}}\) of G, in other words the stabilizer of \(V_i\) ) the vector subspace \(g(V_i)\) is a bijection from I to the set B of vector subspaces of \(k^n\) that are \(\neq \{0\}\) and \(\neq k^n\) , compatible with the action of G.
\item
  If E is a vector space, a flag of E is a set of vector subspaces of E that is totally ordered by inclusion. Show that elements \(a_{1}, \ldots, a_{k}\) of I belong to the same facet of I if and only if \(\{j(a_{1}), \ldots, j(a_{k})\}\) is a flag of \(k^{n}\) .
\item
  Show that G acts doubly-transitively on the set of 1-dimensional vector subspaces of \(k^{n}\) . Suppose that \(n \geqslant 2\) and let \(N_{1}\) be the subgroup of G generated by the element that interchanges \(e_{1}\) and \(e_{2}\) and leaves fixed the other \(e_{i}\) . Show that \((\mathrm{G}, \mathrm{G}^{(1)}, \mathrm{N}_{1})\) is a Tits system whose Weyl group is of order 2.
\item
  Assume that \(k\) is commutative. Set
\end{enumerate}

\[
\mathrm{G} ^ {\prime} = \mathbf {S L} (n, k), \quad \mathrm{B} ^ {\prime} = \mathrm{G} ^ {\prime} \cap \mathrm{B}, \quad \mathrm{N} ^ {\prime} = \mathrm{G} ^ {\prime} \cap \mathrm{N} \text {and} \mathrm{T} ^ {\prime} = \mathrm{N} ^ {\prime} \cap \mathrm{B} ^ {\prime} = \mathrm{T} \cap \mathrm{G} ^ {\prime}.
\]

Show that \(N^{\prime}/T^{\prime}\) can be identified with \(S_{n}\) and that \((G^{\prime}, B^{\prime}, N^{\prime}, S)\) is a Tits system (argue as in no. 2).

\begin{enumerate}
\def\labelenumi{\arabic{enumi})}
\setcounter{enumi}{18}
\item
  Let \(k\) be a commutative field of characteristic \(\neq 2\) and let \(\mathbf{Q}\) be the quadratic form \(x_{1}x_{3} + x_{2}^{2}\) on \(k^{3}\) . Show that the group \(\mathbf{SO}(\mathbf{Q})\) acts doubly transitively on the set of isotropic lines in \(k^{3}\) . Compare the corresponding Tits system (Exerc. 7) with that obtained in no. 2 for \(n = 2\) (cf.~Algebra Chap. IX, §9, Exerc. 15).
\item
  We use the notation of no. 2, and assume that k is commutative. Consider the following cases:
\end{enumerate}

\((B_{l})\ n=2l+1\ (with\ l\geqslant1),\ k\ is\ of\ characteristic\neq2\ and\ k^{n}\ is\ equipped\) with the quadratic form \(Q=x_{1}x_{n}+\cdots+x_{l}x_{l+2}+x_{l+1}^{2};\)

\((C_{l}) \ n = 2l \ (\text{with } l \geqslant 1) \ \text{and} \ k^{n} \ \text{is equipped with the alternating form } \Phi \ \text{such that } \Phi(e_{i}, e_{j}) = 0 \ \text{for} \ 1 \leqslant i \leqslant n, \ 1 \leqslant j \leqslant n \ \text{and} \ i \leqslant j, \ \text{except when} \ i \leqslant l, \ j = n + 1 - i, \ \text{when} \Phi(e_{i}, e_{j}) = 1;\)

\((D_{l})\ n = 2l \text{ (with } l \geqslant 2), k \text{ is of characteristic } \neq 2 \text{ and } k^{n} \text{ is equipped with the quadratic form } Q = x_{1}x_{n} + \cdots + x_{l}x_{l+1}.\)

Denote by \(G_{1}\) the special orthogonal group \(\mathbf{SO}(Q)\) in cases \((B_{l})\) and \((D_{l})\) , and the symplectic group \(\mathbf{Sp}(\varPhi)\) in case \((C_{l})\) . Put

\[
\mathrm{B} _ {1} = \mathrm{G} _ {1} \cap \mathrm{B}, \quad \mathrm{N} _ {1} = \mathrm{G} _ {1} \cap \mathrm{N} \text {and} \mathrm{T} _ {1} = \mathrm{G} _ {1} \cap \mathrm{T} = \mathrm{B} _ {1} \cap \mathrm{N} _ {1}.
\]

\begin{enumerate}
\def\labelenumi{\alph{enumi})}
\tightlist
\item
  Show that the action of \(N_{1}\) on the set of lines \(ke_{i}\) defines a homomorphism from \(N_{1}\) , with kernel \(T_{1}\) , onto the subgroup \(W_{1}\) of \(S_{n}\) , allowing us to identify \(W_{1}\) with \(N_{1}/T_{1}\) . Show that \(W_{1}\) is the subgroup of \(S_{n}\) generated by
\end{enumerate}

the \(\sigma_{j} = s_{j}s_{n - j}\) for \(1\leqslant j <   l\) and by \(\sigma_l = s_{l - 1}s_ls_{l - 1}\) in case \((\mathrm{B}_l)\) ;

the \(\sigma_{j} = s_{j}s_{n - j}\) for \(1\leqslant j <   l\) and by \(\sigma_l = s_l\) in case \((C_l)\) ;

\[
\text { the } \sigma_ {j} = s _ {j} s _ {n - j} \text { for } 1 \leqslant j <   l \text { and by } \sigma_ {l} = s _ {l - 1} s _ {l} s _ {l - 1} s _ {l} s _ {l + 1} s _ {l} \text { in case } (D _ {l}).
\]

\begin{enumerate}
\def\labelenumi{\alph{enumi})}
\setcounter{enumi}{1}
\tightlist
\item
  Let \(S_{1}\) be the set of \(\sigma_{j}\) for \(1 \leqslant j \leqslant l\) . Show that \((\mathrm{G}_{1}, \mathrm{B}_{1}, \mathrm{N}_{1}, \mathrm{S}_{1})\) is a Tits system (to prove that the subgroup H of \(G_{1}\) generated by \(B_{1}\) and \(N_{1}\) is the whole of \(G_{1}\) , argue as in Algebra, Chap. II, §10, no. 13), remarking that H contains the lower triangular subgroup of \(G_{1}\) and that, for any \(\xi_{i} \in k\) ( \(2 \leqslant i \leqslant n\) ), there exists a matrix \(b = (b_{ij}) \in B_{1}\) such that \(b_{11} = 1\) , \(b_{1i} = \xi_{i}\) for \(2 \leqslant i \leqslant n - 1\) and \(b_{1n} = \xi_{n}\) in case \((\mathrm{C}_{l})\) , and \(b_{1n} = 0\) in cases \((\mathrm{B}_{l})\) and \((\mathrm{D}_{l})\) . Then argue as in no. 2, introducing the subgroups \(G_{1,j} = G_{1} \cap (G_{j}, G_{n-j})\) for \(1 \leqslant j < l\) and the subgroup \(G_{1,l}\) of elements of \(G_{1}\) leaving fixed
\end{enumerate}

the \(e_i\) for \(i \neq l, l + 1, l + 2\) and the subspace generated by \(e_l, e_{l+1}\) and \(e_{l+2}\) in case \((B_l)\) ;

the \(e_{i}\) for \(i \neq l, l + 1\) and the plane generated by \(e_{l}\) and \(e_{l+1}\) in case \((\mathrm{C}_{l})\) ; the \(e_{i}\) for i \textless{} l - 1 or \(i > l + 2\) and the two planes generated by \(e_{l-1}\) and \(e_{l+1}\) and by \(e_{l}\) and \(e_{l+2}\) in case \((\mathrm{D}_{l})\) .

Show that \(G_{1,j}\) can be identified, respectively, with \(\mathbf{GL}(2,k)\) , \(\mathbf{SL}(2,k)\) or the special orthogonal group of Exerc. 19).

\(^{*}c)\) Show that the Coxeter graph of the group \(W_{1}\) is of type \((B_{l})\) , \((C_{l})\) or \((D_{l})\) in the three cases (Chap VI, §4, no. 1).

\(d\) ) Show that, for any subset X of \(\mathbf{S}_1\) , the subgroup \(\mathbf{G}_{1\mathbf{X}}\) consists of the \(g\in \mathbf{G}_1\) such that \(g(\mathrm{V}_i) = \mathrm{V}_i\) for all \(i\) such that \(\sigma_{i}\notin \mathbf{X}\) , except in case \((\mathrm{D}_l)\) where the same assertion is valid provided that \(\mathrm{V}_{r - 1}\) denotes the subspace generated by \(e_1,\dots ,e_{r - 1}\) and \(e_{r + 1}\) . Deduce, as in Exerc. 18 b), that there exists a bijection \(j\) from the building associated to \((\mathrm{G}_1,B_1)\) to the set of totally isotropic subspaces \(\neq 0\) in cases \((\mathrm{B}_l)\) and \((\mathrm{C}_l)\) , and to the set of totally isotropic subspaces of dimension \(\neq 0\) and \(\neq r - 1\) in case \((\mathrm{D}_l)\) . Show that points \(a_1,\ldots ,a_k\) of I belong to the same facet if and only if \(\{j(a_1),\dots ,j(a_k)\}\) is a flag.

\begin{enumerate}
\def\labelenumi{\arabic{enumi})}
\setcounter{enumi}{20}
\tightlist
\item
  Let A be a discrete valuation ring (Commutative Algebra, Chap. VI, §3, no. 6), m its maximal ideal, γ a generator of m and K the field of fractions of A. Let G be the group SL(2,K), B the subgroup of G consisting of the matrices \(\begin{pmatrix} a & b \\ c & d \end{pmatrix}\) such that \(a, b, d \in A\) and \(c \in m\) (with ad - bc = 1) and N the subgroup consisting of the matrices belonging to G and having only one non-zero entry in each row and column.
\end{enumerate}

\begin{enumerate}
\def\labelenumi{\alph{enumi})}
\item
  Show that \(\mathrm{T} = \mathrm{B} \cap \mathrm{N}\) is normal in \(\mathbb{N}\) and that \(\mathrm{W} = \mathrm{N} / \mathrm{T}\) is an infinite dihedral group generated by the classes \(s\) and \(s'\) of the matrices \(\left( \begin{array}{cc}0 & 1\\ -1 & 0 \end{array} \right)\) and \(\left( \begin{array}{cc}0 & \gamma \\ -\gamma^{-1} & 0 \end{array} \right)\) , respectively.
\item
  Show that \((G, B, N, S)\) (with \(S = \{s, s'\}\) ) is a Tits system.
\item
  Let \(\mathbf{H} = \mathbf{SL}(2, \mathbf{A})\) be the subgroup of G consisting of the matrices with coefficients in A. Show that \((\mathbf{H}, \mathbf{B} \cap \mathbf{H}, \mathbf{N}, \{s\})\) is a Tits system. Compare with Exerc. 18 e).
\item
  Let \(\hat{\mathbf{A}}\) be the completion of \(\mathbf{A}\) and let \(\hat{\mathbf{G}},\hat{\mathbf{B}},\hat{\mathbf{N}},\hat{\mathbf{T}}\) be the groups defined as above but replacing \(\mathbf{A}\) by \(\hat{\mathbf{A}}\) . Show that the injection of \(\mathbf{G}\) into \(\hat{\mathbf{G}}\) defines an isomorphism from the building \(\mathbf{I}\) associated to \((\mathbf{G},\mathbf{B})\) (Exerc. 10) to the building \(\hat{\mathbf{I}}\) associated to \((\hat{\mathbf{G}},\hat{\mathbf{B}})\) . Let \((\mathbf{I},\mathfrak{U})\) (resp. \((\hat{\mathbf{I}},\hat{\mathfrak{U}})\) ) be the structured building associated to \((\mathbf{G},\mathbf{B},\mathbf{N})\) (resp. \((\hat{\mathbf{G}},\hat{\mathbf{B}},\hat{\mathbf{N}})\) ) (Exerc. 12): show that \(j(\mathfrak{U})\subset \hat{\mathfrak{U}}\) , but that \(j(\mathfrak{U})\neq \hat{\mathfrak{U}}\) if \(\mathbf{A}\neq \hat{\mathbf{A}}\) (remark that the apartments \(\hat{\mathfrak{U}}\) (resp. \(\mathfrak{U}\) ) correspond bijectively to the conjugates of \(\mathbf{T}\) (resp. \(\hat{\mathbf{T}}\) ) by \(\mathbf{G}\) (resp. \(\hat{\mathbf{G}}\) ).
\end{enumerate}

\begin{enumerate}
\def\labelenumi{\arabic{enumi})}
\setcounter{enumi}{21}
\tightlist
\item
  Let \(G\) be a group and \(B\) a subgroup of \(G\) .
\end{enumerate}

\begin{enumerate}
\def\labelenumi{\alph{enumi})}
\tightlist
\item
  Show that the following conditions are equivalent:
\end{enumerate}

\begin{enumerate}
\def\labelenumi{(\roman{enumi})}
\item
  \(B \cap gBg^{-1}\) is of finite index in \(B\) for all \(g \in G\) ;
\item
  every double coset BgB with respect to B is a finite union of left cosets with respect to B.
\end{enumerate}

More precisely, show that, for all \(g \in G\) , the index \(q_{g}\) of \(B \cap gBg^{-1}\) in B is equal to the number of left cosets with respect to B contained in the double coset BgB. Show that \(q_{gh} \leqslant q_{g}q_{h}\) for all \(g, h \in G\) .

Assume from now on that conditions (i) and (ii) are satisfied and denote by k a commutative ring. For \(t \in G/B\) (resp. \(t \in B \backslash G/B\) ), denote by \(a_{t}\) the map from G to k defined by \(a_{t}(g) = 0\) if \(g \notin t\) and \(a_{t}(g) = 1\) if \(g \in t\) . Let L (resp. H) be the k-module generated by the \(a_{t}\) for \(t \in G/B\) (resp. \(t \in B \backslash G/B\) ).

\begin{enumerate}
\def\labelenumi{\alph{enumi})}
\setcounter{enumi}{1}
\item
  Show that there exists a unique linear form \(\mu\) on \(\mathbf{L}\) such that \(\mu(a_t) = 1\) for all \(t \in \mathbf{G} / \mathbf{B}\) .
\item
  Let \(\varphi \in \mathbf{L}\) and \(\psi \in \mathbf{H}\) . Show that, for all \(x \in \mathbf{G}\) , the map
\end{enumerate}

\[
\theta_ {x}: y \mapsto \varphi (y) \psi (y ^ {- 1} x)
\]

belongs to L and that the map \(\varphi * \psi : x \mapsto \mu(\theta_x)\) belongs to L. Moreover, if \(\varphi \in H\) then \(\varphi * \psi \in H\) . Show that the map \((\varphi, \psi) \mapsto \varphi * \psi\) makes H into an algebra over \(k\) , having \(a_B\) as its unit element, and makes L into a right H-module.

The algebra \(\mathbf{H}\) is called the Hecke algebra of \(\mathbf{G}\) with respect to \(\mathbf{H}\) and is denoted by \(\mathrm{H}_k(\mathbf{G},\mathbf{B})\) .

\begin{enumerate}
\def\labelenumi{\alph{enumi})}
\setcounter{enumi}{3}
\tightlist
\item
  Show that, for \(t, t' \in B \backslash G / B\) ,
\end{enumerate}

\[
a _ {t} * a _ {t ^ {\prime}} = \sum_ {t ^ {\prime \prime}} m (t, t ^ {\prime}; t ^ {\prime \prime}) a _ {t ^ {\prime \prime}},
\]

where \(m(t, t'; t'')\) is the number of cosets with respect to B contained in \(t \cap gt'^{-1}\) for all \(g \in t''\) .

\begin{enumerate}
\def\labelenumi{\alph{enumi})}
\setcounter{enumi}{4}
\tightlist
\item
  Let G act on L by left translations. Show that the action of H on L defines an isomorphism from H to the commutant of the linear representation of G on L thus obtained.
\end{enumerate}

\(^{*}f\) ) Assume that G is finite and that the characteristic of k does not divide the order of G. Show that \(H_{k}(G,B)\) is absolutely semi-simple over k (Algebra, Chap VIII, §7, no. 5) (use Maschke's Theorem (Chap. V, Appendix) and Prop. 3 of Algebra, Chap. VIII, §5, no. 1).

\begin{enumerate}
\def\labelenumi{\alph{enumi})}
\setcounter{enumi}{6}
\tightlist
\item
  Assume that G is a topological group and that B is a compact open subgroup of G. Show that conditions (i) and (ii) are satisfied and that, when k = R or C, the product \(\varphi * \psi\) is simply the convolution product relative to the right Haar measure on G, normalised by the condition \(\mu(B) = 1\) (cf.~Integration, Chap. VIII, §4, no. 5).
\end{enumerate}
